% Options for packages loaded elsewhere
\PassOptionsToPackage{unicode}{hyperref}
\PassOptionsToPackage{hyphens}{url}
\documentclass[
  a4paper,
]{article}
\usepackage{xcolor}
\usepackage[margin=25mm]{geometry}
\usepackage{amsmath,amssymb}
\setcounter{secnumdepth}{-\maxdimen} % remove section numbering
\usepackage{iftex}
\ifPDFTeX
  \usepackage[T1]{fontenc}
  \usepackage[utf8]{inputenc}
  \usepackage{textcomp} % provide euro and other symbols
\else % if luatex or xetex
  \usepackage{unicode-math} % this also loads fontspec
  \defaultfontfeatures{Scale=MatchLowercase}
  \defaultfontfeatures[\rmfamily]{Ligatures=TeX,Scale=1}
\fi
\usepackage{lmodern}
\ifPDFTeX\else
  % xetex/luatex font selection
\fi
% Use upquote if available, for straight quotes in verbatim environments
\IfFileExists{upquote.sty}{\usepackage{upquote}}{}
\IfFileExists{microtype.sty}{% use microtype if available
  \usepackage[]{microtype}
  \UseMicrotypeSet[protrusion]{basicmath} % disable protrusion for tt fonts
}{}
\makeatletter
\@ifundefined{KOMAClassName}{% if non-KOMA class
  \IfFileExists{parskip.sty}{%
    \usepackage{parskip}
  }{% else
    \setlength{\parindent}{0pt}
    \setlength{\parskip}{6pt plus 2pt minus 1pt}}
}{% if KOMA class
  \KOMAoptions{parskip=half}}
\makeatother
\usepackage{color}
\usepackage{fancyvrb}
\newcommand{\VerbBar}{|}
\newcommand{\VERB}{\Verb[commandchars=\\\{\}]}
\DefineVerbatimEnvironment{Highlighting}{Verbatim}{commandchars=\\\{\}}
% Add ',fontsize=\small' for more characters per line
\newenvironment{Shaded}{}{}
\newcommand{\AlertTok}[1]{\textcolor[rgb]{1.00,0.00,0.00}{\textbf{#1}}}
\newcommand{\AnnotationTok}[1]{\textcolor[rgb]{0.38,0.63,0.69}{\textbf{\textit{#1}}}}
\newcommand{\AttributeTok}[1]{\textcolor[rgb]{0.49,0.56,0.16}{#1}}
\newcommand{\BaseNTok}[1]{\textcolor[rgb]{0.25,0.63,0.44}{#1}}
\newcommand{\BuiltInTok}[1]{\textcolor[rgb]{0.00,0.50,0.00}{#1}}
\newcommand{\CharTok}[1]{\textcolor[rgb]{0.25,0.44,0.63}{#1}}
\newcommand{\CommentTok}[1]{\textcolor[rgb]{0.38,0.63,0.69}{\textit{#1}}}
\newcommand{\CommentVarTok}[1]{\textcolor[rgb]{0.38,0.63,0.69}{\textbf{\textit{#1}}}}
\newcommand{\ConstantTok}[1]{\textcolor[rgb]{0.53,0.00,0.00}{#1}}
\newcommand{\ControlFlowTok}[1]{\textcolor[rgb]{0.00,0.44,0.13}{\textbf{#1}}}
\newcommand{\DataTypeTok}[1]{\textcolor[rgb]{0.56,0.13,0.00}{#1}}
\newcommand{\DecValTok}[1]{\textcolor[rgb]{0.25,0.63,0.44}{#1}}
\newcommand{\DocumentationTok}[1]{\textcolor[rgb]{0.73,0.13,0.13}{\textit{#1}}}
\newcommand{\ErrorTok}[1]{\textcolor[rgb]{1.00,0.00,0.00}{\textbf{#1}}}
\newcommand{\ExtensionTok}[1]{#1}
\newcommand{\FloatTok}[1]{\textcolor[rgb]{0.25,0.63,0.44}{#1}}
\newcommand{\FunctionTok}[1]{\textcolor[rgb]{0.02,0.16,0.49}{#1}}
\newcommand{\ImportTok}[1]{\textcolor[rgb]{0.00,0.50,0.00}{\textbf{#1}}}
\newcommand{\InformationTok}[1]{\textcolor[rgb]{0.38,0.63,0.69}{\textbf{\textit{#1}}}}
\newcommand{\KeywordTok}[1]{\textcolor[rgb]{0.00,0.44,0.13}{\textbf{#1}}}
\newcommand{\NormalTok}[1]{#1}
\newcommand{\OperatorTok}[1]{\textcolor[rgb]{0.40,0.40,0.40}{#1}}
\newcommand{\OtherTok}[1]{\textcolor[rgb]{0.00,0.44,0.13}{#1}}
\newcommand{\PreprocessorTok}[1]{\textcolor[rgb]{0.74,0.48,0.00}{#1}}
\newcommand{\RegionMarkerTok}[1]{#1}
\newcommand{\SpecialCharTok}[1]{\textcolor[rgb]{0.25,0.44,0.63}{#1}}
\newcommand{\SpecialStringTok}[1]{\textcolor[rgb]{0.73,0.40,0.53}{#1}}
\newcommand{\StringTok}[1]{\textcolor[rgb]{0.25,0.44,0.63}{#1}}
\newcommand{\VariableTok}[1]{\textcolor[rgb]{0.10,0.09,0.49}{#1}}
\newcommand{\VerbatimStringTok}[1]{\textcolor[rgb]{0.25,0.44,0.63}{#1}}
\newcommand{\WarningTok}[1]{\textcolor[rgb]{0.38,0.63,0.69}{\textbf{\textit{#1}}}}
\usepackage{longtable,booktabs,array}
\newcounter{none} % for unnumbered tables
\usepackage{calc} % for calculating minipage widths
% Correct order of tables after \paragraph or \subparagraph
\usepackage{etoolbox}
\makeatletter
\patchcmd\longtable{\par}{\if@noskipsec\mbox{}\fi\par}{}{}
\makeatother
% Allow footnotes in longtable head/foot
\IfFileExists{footnotehyper.sty}{\usepackage{footnotehyper}}{\usepackage{footnote}}
\makesavenoteenv{longtable}
\usepackage{graphicx}
\makeatletter
\newsavebox\pandoc@box
\newcommand*\pandocbounded[1]{% scales image to fit in text height/width
  \sbox\pandoc@box{#1}%
  \Gscale@div\@tempa{\textheight}{\dimexpr\ht\pandoc@box+\dp\pandoc@box\relax}%
  \Gscale@div\@tempb{\linewidth}{\wd\pandoc@box}%
  \ifdim\@tempb\p@<\@tempa\p@\let\@tempa\@tempb\fi% select the smaller of both
  \ifdim\@tempa\p@<\p@\scalebox{\@tempa}{\usebox\pandoc@box}%
  \else\usebox{\pandoc@box}%
  \fi%
}
% Set default figure placement to htbp
\def\fps@figure{htbp}
\makeatother
\setlength{\emergencystretch}{3em} % prevent overfull lines
\providecommand{\tightlist}{%
  \setlength{\itemsep}{0pt}\setlength{\parskip}{0pt}}
\usepackage{bookmark}
\IfFileExists{xurl.sty}{\usepackage{xurl}}{} % add URL line breaks if available
\urlstyle{same}
\hypersetup{
  pdftitle={The Sixth Thought Experiment: Can Two Kinds of Long-Range Interaction Be Internalized into a Single Closed State Map? - Reconstructing Gravity{,} Coulomb{,} GW-Quadrupole and EM-Dipole Radiation of a Charged Quasi-Circular Binary with Eight Persistent States{,} and Verifying the Transferability of the Same Transition over Five Conditions -},
  pdfauthor={Noriaki Kihara (ORCID 0009-0004-6753-4020)},
  hidelinks,
  pdfcreator={LaTeX via pandoc}}

\title{The Sixth Thought Experiment: Can Two Kinds of Long-Range
Interaction Be Internalized into a Single Closed State Map? -
Reconstructing Gravity, Coulomb, GW-Quadrupole and EM-Dipole Radiation
of a Charged Quasi-Circular Binary with Eight Persistent States, and
Verifying the Transferability of the Same Transition over Five
Conditions -}
\author{Noriaki Kihara (ORCID 0009-0004-6753-4020)}
\date{2026-09-26 / v1.1 / Version DOI 10.5281/zenodo.22974632}

\sloppy
\emergencystretch=3em
\begin{document}
\maketitle

\textbf{Author:} Noriaki Kihara\\
\textbf{ORCID:} 0009-0004-6753-4020\\
\textbf{Version DOI:} 10.5281/zenodo.22974632\\
\textbf{Concept DOI:} 10.5281/zenodo.22974631\\
\textbf{Version:} v1.1\\
\textbf{Date:} 2026-09-26

\textbf{Keywords:} charged binary, state closure, synchronous state map,
Coulomb interaction, gravitational-wave radiation, electric-dipole
radiation, radiation reaction, persistent logical state, discrete
dynamics, reproducibility

\begin{center}\rule{0.5\linewidth}{0.5pt}\end{center}

\section{Abstract}\label{abstract}

The fifth thought experiment {[}5{]} self-audited, from the viewpoint of
information flow, the gravitational-orbit generator that the fourth
thought experiment {[}4{]} had displayed as two states, and confirmed
that in the present representation the information held across macro
steps must be made explicit as

\[
Z_6=(U,P,E,H,Q,N).
\]

Furthermore, the RK4 intermediate quantities and the processing phase
were placed inside the state as machine/work states, and the generator
was reconstructed into an OLD-state-only synchronous update, with no
external physical parameters, a non-feedback readout, and a single
\texttt{transition(z)} as the only generating path.

In this paper those audit rules are applied for the first time to
\textbf{a system extended from one kind of interaction to two kinds}.
The object is a charged two-body system restricted to \(G=M=c=1\),
non-spinning point particles, quasi-circular orbits and a leading-order
adiabatic inspiral. The conservative motion uses Newtonian gravity and
the Coulomb interaction, and the dissipation uses the leading
gravitational quadrupole radiation and the leading electric-dipole
radiation. On the independent analytic reference side, the charge
dependence can be organized into the two relational quantities

\[
C=1-\lambda_A\lambda_B,
\qquad
D=(\lambda_A-\lambda_B)^2,
\]

where \(\lambda_A=q_A/m_A\) and \(\lambda_B=q_B/m_B\). \(C\) enters the
conservative force and the orbital angular velocity, and \(D\) decides
the presence and the strength of the leading electric-dipole radiation.

Adding \(C,D\) to the six persistent states of the fifth thought
experiment gives

\[
\boxed{
Z_8=(U,P,E,H,Q,N,C,D).
}
\]

In the implementation, a full microstate of 41 real components is used,
consisting of these eight persistent logical states, the RK4 work
registers and the 11-phase one-hot phase. \texttt{transition(z)}
receives no external physical argument, builds every candidate value
from the OLD full state only, and generates the next state
simultaneously by a one-hot sum of products. Since \(N,C,D\) are held as
identity candidates in all phases, they are preserved as the result of
the same map, not by an external copy.

For the reference case

\[
\lambda_A=+0.30,
\qquad
\lambda_B=-0.30,
\]

that is,

\[
C=1.09,
\qquad
D=0.36,
\qquad
N=\nu=0.25,
\]

the strict eight-state generation from \(r=50\rightarrow20\) completed
in 488,345 macro steps. The differences from the analytic reference were

\[
|\Delta t|=1.46\times10^{-8},
\qquad
|\Delta\phi|=8.44\times10^{-11},
\]

and over the whole stored range

\[
\max|\Delta r|=9.12\times10^{-9},
\]

with

\[
\Delta N=\Delta C=\Delta D=0.
\]

Furthermore, with the same \texttt{transition(z)}, the same 11-phase
machine and the same eight persistent states unchanged, the five
conditions differing only in the initial \(C,D\),

\[
(+0.30,-0.30),
(+0.30,+0.30),
(+0.90,-0.90),
(+0.90,+0.90),
(+1.20,-1.20),
\]

were run. In all five conditions the orbital state generation completed,
and the drift of \(N,C,D\) was 0. In particular, for equal-sign
equal-ratio cases \(D=0\), so that only the leading electric-dipole term
adopted in this paper vanishes, while the Coulomb repulsion effect with
\(C<1\) remains in the conservative motion. The \texttt{n3\_repulsive}
case with \(C=0.19\) required about 3816.7 orbits; \(P,H\) could be
generated to the end, but the time readout through the exponential clock
\(Q\) exceeded the finite-precision range. This fact is not corrected;
it is recorded as it is, as a limit of the domain and of the numerical
representation.

What this paper shows is that, for the restricted leading-order
quasi-circular charged binary, the conservative effects of gravity and
Coulomb and the dissipative effects of GW quadrupole and EM dipole
radiation \textbf{can be rearranged into a single synchronous state map
with no external charge parameters and no branching by force type}. This
does not mean a fundamental unification of gravity and electromagnetism,
a derivation of the Einstein--Maxwell equations, the minimality of eight
states, the ontological fundamentality of \(C,D\), or universality for
general charged binaries.

\begin{center}\rule{0.5\linewidth}{0.5pt}\end{center}

\section{0. Position of this paper}\label{position-of-this-paper}

This paper is the sixth thought experiment under the research project
design {[}0{]}.

The first thought experiment {[}1{]} examined, without placing a
background spacetime or a metric first, the minimal construction that
distinguishes inertial, uniformly accelerated and rotating motion, from
two degrees of freedom of unfixed meaning and the conservation of an
area readout.

The second thought experiment {[}2{]} applied a fixed action repeatedly
to two values and constructed Coulomb-type inverse-square motion from a
product readout and an area clock.

The third thought experiment {[}3{]} reversed the direction and treated
the inverse problem of reading out the mass ratio, the relative distance
and the relative time from known two-body orbits.

The fourth thought experiment {[}4{]} brought the gravitational orbits
computed externally in the third thought experiment back to an internal
sequential generation by the state-dependent action

\[
X_{n+1}=S_nX_n.
\]

The fifth thought experiment {[}5{]} self-audited that fourth thought
experiment. What became important there was the point that

\begin{center}
\fbox{\parbox{0.90\linewidth}{\centering even if a numerical result agrees with a known orbit, the system is not necessarily closed by the explicit states alone.}}
\end{center}

Since information held across macro steps remained outside the apparent
two states \(X=(a,b)^T\) of the fourth thought experiment, it was made
explicit as the six persistent logical states

\[
Z_6=(U,P,E,H,Q,N),
\]

and the intermediate RK4 values and the 11-phase one-hot phase were also
placed inside the full state as machine/work states.

The sixth thought experiment does not change these audit rules
themselves. The new question is

\[
\boxed{
\begin{array}{c}
\text{when the interaction is increased from one kind, gravity, to two kinds, gravity + Coulomb,}\\
\text{and the dissipation too is increased to the two channels GW quadrupole + EM dipole,}\\
\text{can the system still be closed as a single synchronous state map without external physical information?}
\end{array}
}
\]

The subject of this paper is therefore not merely ``computing charged
orbits''. With the known charged binary placed on the answer-checking
side, it is to audit \textbf{what must be held internally as persistent
states so that the same update rule transfers to several conditions}.

\begin{center}\rule{0.5\linewidth}{0.5pt}\end{center}

\section{1. Motivation}\label{motivation}

\subsection{1.1 From one kind of interaction to
two}\label{from-one-kind-of-interaction-to-two}

In a gravity-only binary, there is one kind of central interaction. When
the Coulomb interaction is added, on the conservative side the same
\(1/r^2\) type changes between attraction and repulsion according to the
sign. On the dissipative side, moreover, a leading electric-dipole
radiation appears, distinct from the gravitational quadrupole radiation
{[}6--11{]}.

The charged binary is therefore the minimal next stage that tests

\begin{center}
\fbox{\parbox{0.90\linewidth}{\centering multiplication of conservative interactions$+$multiplication of dissipation channels}}
\end{center}

at the same time.

On the other hand, going to the general Einstein--Maxwell two-body
problem makes the object of verification rapidly complex. This paper
deliberately restricts itself to

\begin{itemize}
\tightlist
\item
  non-spinning point particles,
\item
  quasi-circular orbits,
\item
  leading-order conservative dynamics,
\item
  leading GW quadrupole radiation,
\item
  leading EM electric-dipole radiation,
\item
  adiabatic energy balance,
\item
  \(G=M=c=1\).
\end{itemize}

Within this range an independent analytic reference can be stated
explicitly, and the direction of comparison with the state generator can
be fixed to the one direction

\[
\boxed{
\text{state generator output}
\longrightarrow
\text{independent analytic reference}.
}
\]

\subsection{1.2 Not supplying the charge parameters from outside at
every
step}\label{not-supplying-the-charge-parameters-from-outside-at-every-step}

If a charged binary is implemented in the ordinary way, at each step

\[
\lambda_A,
\qquad
\lambda_B
\]

can be referred to as external parameters. Seen from the audit rules of
the fifth thought experiment {[}5{]}, however, ``the physical
information that determines the next state'' then remains outside the
state.

This paper therefore first audits the combinations that actually appear
in the analytic formulas. For the present restricted model it turns out
that, instead of holding the individual \(\lambda_A,\lambda_B\)
directly,

\[
\boxed{
C=1-\lambda_A\lambda_B
}
\]

and

\[
\boxed{
D=(\lambda_A-\lambda_B)^2
}
\]

suffice.

\(C\) enters the conservative force and the orbital frequency, and \(D\)
enters the leading electric-dipole radiation. If these two can be
internalized into the state, \texttt{transition(z)} needs neither the
individual labels of bodies A/B, nor an external attraction/repulsion
branch, nor a gravity/electromagnetism force-type branch.

\subsection{1.3 Looking at transferability, not agreement in one
condition}\label{looking-at-transferability-not-agreement-in-one-condition}

Even if one condition agrees with the analytic reference, the
possibility remains that the coefficients were made for that case only.
This paper therefore keeps the same \texttt{transition(z)} fixed and
changes only the \(C,D\) put into the initial state.

In particular, by including

\begin{itemize}
\tightlist
\item
  opposite-sign charges: \(C>1\), \(D>0\),
\item
  equal-sign equal charge-to-mass ratios: \(C<1\), \(D=0\),
\end{itemize}

the strength of the conservative force and the presence of dipole
radiation are changed at the same time.

The success criterion is

\begin{center}
\fbox{\parbox{0.90\linewidth}{\centering that the same transition handles different conditions by the difference of state values alone$.$}}
\end{center}

\begin{center}\rule{0.5\linewidth}{0.5pt}\end{center}

\section{2. The question and the scope of the
claims}\label{the-question-and-the-scope-of-the-claims}

The question verified in this paper is divided into the following four
stages.

\begin{enumerate}
\def\labelenumi{\arabic{enumi}.}
\item
  \textbf{Analytic closure}\\
  Can the charge dependence of the present leading-order quasi-circular
  charged binary be organized into \(C,D\)?
\item
  \textbf{State closure}\\
  With \(C,D\) added to the six persistent states of the fifth thought
  experiment, \[
  Z_8=(U,P,E,H,Q,N,C,D),
  \] can the physical case values needed to determine the next state be
  placed inside the state?
\item
  \textbf{Single map}\\
  Can the full state, including the RK4 intermediate quantities, the
  processing phase and the conserved relational quantities, be updated
  by a single \texttt{transition(z)} that takes no external physical
  argument?
\item
  \textbf{Transferability}\\
  Can the same \texttt{transition(z)}, unchanged, be transferred to five
  different conditions by changing only the \(C,D\) of the initial
  state?
\end{enumerate}

The affirmative conclusions of this paper are stated \textbf{only for
this restricted model and this range of implementation}.

\begin{center}\rule{0.5\linewidth}{0.5pt}\end{center}

\section{3. Independent analytic
reference}\label{independent-analytic-reference}

\subsection{3.1 Approximations and
units}\label{approximations-and-units}

The analytic reference is constructed independently of the state
generator that follows. The adopted conditions are

\[
G=M=c=1,
\qquad
\epsilon_0=\frac{1}{4\pi},
\]

\[
M=m_A+m_B=1.
\]

The object is non-spinning point particles in an adiabatic
quasi-circular inspiral; the conservative dynamics is the leading order
of Newtonian gravity + Coulomb, and the dissipation is the leading
mass-quadrupole gravitational radiation and the leading electric-dipole
electromagnetic radiation. Peters--Mathews {[}6,7{]} is the basis of
gravitational quadrupole radiation and orbital decay, and as prior work
treating both GW and EM in charged binaries, Liu et al.~{[}8{]},
Benavides-Gallego and Han {[}9{]}, Zhang et al.~{[}10{]} and Verma et
al.~{[}11{]} are referred to.

In the notation of Zhang et al.~{[}10{]} a quantity corresponding to
\(1-\lambda_A\lambda_B\) is used; in this paper it is written \(C\) to
avoid a clash with the state vector \(Z_8\).

\subsection{3.2 Dimensionless mass and charge
quantities}\label{dimensionless-mass-and-charge-quantities}

The charge-to-mass ratios are

\[
\lambda_A=\frac{q_A}{m_A},
\qquad
\lambda_B=\frac{q_B}{m_B}.
\]

The reduced mass and the symmetric mass ratio are

\[
\mu=\frac{m_A m_B}{M},
\qquad
\nu=\frac{m_A m_B}{M^2}.
\]

For \(M=1\),

\[
\mu=\nu.
\]

Since this experiment has equal masses,

\[
\mu=\nu=\frac14.
\]

\subsection{\texorpdfstring{3.3 The relational quantity \(C\) that
collects the conservative
motion}{3.3 The relational quantity C that collects the conservative motion}}\label{the-relational-quantity-c-that-collects-the-conservative-motion}

Collecting the leading-order effective central interaction of gravity
and Coulomb, put

\[
\boxed{
C=1-\lambda_A\lambda_B.
}
\]

For quasi-circular orbits,

\[
\boxed{
\omega^2=\frac{C}{r^3},
}
\]

and the binding energy is

\[
\boxed{
E_{\rm orb}=-\frac{\mu C}{2r}.
}
\]

If \(\lambda_A\lambda_B<0\) then \(C>1\), and Coulomb acts as an
attraction in the same direction as gravity. If \(\lambda_A\lambda_B>0\)
then \(C<1\), and the Coulomb repulsion weakens the effective central
attraction.

The real quasi-circular reference of this paper requires

\[
\boxed{C>0}.
\]

\(C\le0\) is therefore outside the domain of the present analytic
reference.

\subsection{3.4 GW quadrupole radiation}\label{gw-quadrupole-radiation}

The leading mass-quadrupole power is

\[
P_{\rm GW}
=
\frac{32}{5}\mu^2 r^4\omega^6.
\]

Substituting \(\omega^2=C/r^3\),

\[
\boxed{
P_{\rm GW}
=
\frac{32}{5}
\frac{\mu^2C^3}{r^5}
}
\]

{[}6,7{]}.

\subsection{3.5 EM electric-dipole
radiation}\label{em-electric-dipole-radiation}

The electric dipole moment in the COM frame is, at leading order,

\[
\mathbf d
=
\mu(\lambda_A-\lambda_B)\mathbf r.
\]

Here put

\[
\boxed{
D=(\lambda_A-\lambda_B)^2.
}
\]

The leading dipole power is

\[
\boxed{
P_{\rm EM}
=
\frac23
\frac{\mu^2 D C^2}{r^4}
}
\]

{[}8--11{]}.

From this expression, if

\[
\lambda_A=\lambda_B
\quad\Longrightarrow\quad
D=0,
\]

then \textbf{the leading electric-dipole term adopted in this paper}
vanishes. The Coulomb conservative force itself, however, does not
vanish.

\subsection{3.6 Energy balance and radius
evolution}\label{energy-balance-and-radius-evolution}

The radiative loss is taken as

\[
\frac{dE_{\rm orb}}{dt}
=-(P_{\rm GW}+P_{\rm EM}).
\]

From

\[
\frac{dE_{\rm orb}}{dr}
=
\frac{\mu C}{2r^2},
\]

we obtain

\[
\boxed{
\frac{dr}{dt}
=-\frac{A}{r^2}-\frac{B}{r^3}.
}
\]

Here

\[
\boxed{
A=\frac43\mu D C
}
\]

is the coefficient on the EM dipole side and

\[
\boxed{
B=\frac{64}{5}\mu C^2
}
\]

is the coefficient on the GW quadrupole side.

In the radius evolution, therefore, two dissipative contributions with
different distance dependences,

\[
\boxed{
r^{-2}\ \text{and}\ r^{-3},
}
\]

appear at the same time.

\subsection{\texorpdfstring{3.7 Analytic \(t(r)\) and
\(\phi(r)\)}{3.7 Analytic t(r) and \textbackslash phi(r)}}\label{analytic-tr-and-phir}

Since

\[
\frac{dt}{dr}
=-\frac{r^3}{Ar+B},
\]

for \(A\ne0\), putting

\[
F_t(r)
=
\frac{r^3}{3A}
-\frac{Br^2}{2A^2}
+\frac{B^2r}{A^3}
-\frac{B^3}{A^4}\ln(Ar+B),
\]

we have

\[
\boxed{
t(r)=F_t(r_0)-F_t(r)}.
\]

Also, from

\[
\frac{d\phi}{dt}=\sqrt{\frac{C}{r^3}},
\]

\[
\frac{d\phi}{dr}
=-\frac{\sqrt C\,r^{3/2}}{Ar+B}.
\]

For \(A\ne0\), with

\[
F_\phi(r)
=
\sqrt C
\left[
\frac{2r^{3/2}}{3A}
-\frac{2B\sqrt r}{A^2}
+\frac{2B^{3/2}}{A^{5/2}}
\tan^{-1}\sqrt{\frac{Ar}{B}}
\right],
\]

\[
\boxed{
\phi(r)=F_\phi(r_0)-F_\phi(r)}.
\]

For equal-sign equal ratios, \(D=0\) and hence \(A=0\). In this case the
limit is used directly,

\[
\boxed{
t(r)=\frac{r_0^4-r^4}{4B}}
\]

and

\[
\boxed{
\phi(r)=\frac{2\sqrt C}{5B}
\left(r_0^{5/2}-r^{5/2}\right)
}.
\]

The relative orbit is read out as

\[
\boxed{
x=r\cos\phi,
\qquad
y=r\sin\phi}.
\]

\begin{center}\rule{0.5\linewidth}{0.5pt}\end{center}

\section{4. The analytic reference
case}\label{the-analytic-reference-case}

The reference condition is

\[
m_A=m_B=0.5,
\]

\[
\lambda_A=+0.30,
\qquad
\lambda_B=-0.30.
\]

Then

\[
C=1.09,
\qquad
D=0.36,
\qquad
\mu=\nu=0.25.
\]

The independent analytic reference for \(r=50\rightarrow20\) is

\[
t_f=169032.31760706357,
\]

\[
\phi_f=767.0889994015913,
\]

\[
N_{\rm cyc}=\frac{\phi_f}{2\pi}
=122.086006046.
\]

In the reference case,

\[
\frac{P_{\rm EM}}{P_{\rm GW}}
=
\frac{5}{48}\frac{D}{C}r,
\]

so that at

\[
r_{\rm cross}
=
\frac{48C}{5D}
\simeq29.0667
\]

the EM dipole and the GW quadrupole are comparable. Outside, the EM
dipole is relatively stronger; going inward, the \(r^{-5}\) GW
quadrupole overtakes it.

\begin{center}\rule{0.5\linewidth}{0.5pt}\end{center}

\section{5. Extension from six states to eight
states}\label{extension-from-six-states-to-eight-states}

\subsection{5.1 The six persistent states of the fifth thought
experiment}\label{the-six-persistent-states-of-the-fifth-thought-experiment}

In the fifth thought experiment {[}5{]}, the macro persistent
information needed for the present gravitational-orbit reconstruction
was taken to be

\[
Z_6=(U,P,E,H,Q,N).
\]

Here

\[
U=e^{i\chi},
\qquad
H=e^{i\phi},
\qquad
Q=e^{t/\tau_0},
\qquad
N=\nu.
\]

\(P\) is the state corresponding to the orbital radius / semi-latus
rectum, and \(E\) is the eccentricity state. In the quasi-circular
reference of this paper the update rate of \(E\) is 0, and it is kept at
its initial value 0.

\subsection{5.2 The conserved relational quantities added in the charged
system}\label{the-conserved-relational-quantities-added-in-the-charged-system}

Looking at the present analytic reference from the viewpoint of the
information needed to determine the next state, what is newly needed is

\[
C=1-\lambda_A\lambda_B
\]

and

\[
D=(\lambda_A-\lambda_B)^2.
\]

Accordingly,

\[
\boxed{
Z_8=(U,P,E,H,Q,N,C,D).
}
\]

What matters here is not to assert that \(C,D\) are ``new fundamental
physical degrees of freedom''. They are \textbf{the persistent
relational quantities needed to close the present restricted analytic
representation}.

\subsection{5.3 Table of persistent
states}\label{table-of-persistent-states}

{\def\LTcaptype{none} % do not increment counter
\begin{longtable}[]{@{}
  >{\raggedright\arraybackslash}p{(\linewidth - 6\tabcolsep) * \real{0.2500}}
  >{\raggedright\arraybackslash}p{(\linewidth - 6\tabcolsep) * \real{0.2500}}
  >{\raggedright\arraybackslash}p{(\linewidth - 6\tabcolsep) * \real{0.2500}}
  >{\raggedright\arraybackslash}p{(\linewidth - 6\tabcolsep) * \real{0.2500}}@{}}
\toprule\noalign{}
\begin{minipage}[b]{\linewidth}\raggedright
State
\end{minipage} & \begin{minipage}[b]{\linewidth}\raggedright
Present meaning
\end{minipage} & \begin{minipage}[b]{\linewidth}\raggedright
Kind
\end{minipage} & \begin{minipage}[b]{\linewidth}\raggedright
Update across macro steps
\end{minipage} \\
\midrule\noalign{}
\endhead
\bottomrule\noalign{}
\endlastfoot
\(U\) & \(e^{i\chi}\) & Phase state & Rotation \\
\(P\) & Orbital radius / semi-latus rectum state & Dynamical & Decreases
by radiation reaction \\
\(E\) & Eccentricity state & Dynamical state & \(dE/d\phi=0\) in this
quasi-circular model \\
\(H\) & \(e^{i\phi}\) & Observed azimuth phase & Rotation \\
\(Q\) & \(e^{t/\tau_0}\) & Clock encoding & Multiplicative update \\
\(N\) & \(\nu\) & Conserved relational quantity & \(N'=N\) \\
\(C\) & \(1-\lambda_A\lambda_B\) & Conserved relational quantity &
\(C'=C\) \\
\(D\) & \((\lambda_A-\lambda_B)^2\) & Conserved relational quantity &
\(D'=D\) \\
\end{longtable}
}

The individual \(\lambda_A,\lambda_B\) are not passed to
\texttt{transition(z)}. The difference between attraction and repulsion,
and the presence or absence of the EM dipole, are expressed as
differences of the state values \(C,D\).

\begin{center}\rule{0.5\linewidth}{0.5pt}\end{center}

\section{6. The state generator}\label{the-state-generator}

\subsection{6.1 Local update formulas with the phase as the independent
variable}\label{local-update-formulas-with-the-phase-as-the-independent-variable}

Combining

\[
\frac{dr}{dt}
=-\frac{A}{r^2}-\frac{B}{r^3}
\]

of the analytic reference with

\[
\frac{d\phi}{dt}=\sqrt{\frac{C}{r^3}},
\]

\[
\frac{dr}{d\phi}
=
-\frac{4}{3}\mu D\frac{\sqrt C}{\sqrt r}
-\frac{64}{5}\mu C\frac{\sqrt C}{r\sqrt r}.
\]

For \(M=1\), \(\mu=\nu=N\), so in state variables

\[
\boxed{
\frac{dP}{d\phi}
=
-\frac{4}{3}ND\frac{\sqrt C}{\sqrt P}
-\frac{64}{5}NC\frac{\sqrt C}{P\sqrt P}.
}
\]

Also

\[
\boxed{
\frac{dE}{d\phi}=0
}
\]

and

\[
\boxed{
\frac{dt}{d\phi}
=
\frac{P\sqrt P}{\sqrt C}.
}
\]

The local rate

\[
R(P,E,N,C,D)
=
\left(
\frac{dP}{d\phi},
0,
\frac{dt}{d\phi}
\right)
\]

can thus be evaluated \textbf{from the present states \(P,N,C,D\)
alone}.

\subsection{6.2 Numerical resolution and encoding
constants}\label{numerical-resolution-and-encoding-constants}

One orbit is divided into 4000 macro steps,

\[
h=\frac{2\pi}{4000}.
\]

The clock state is encoded as

\[
Q=e^{t/\tau_0},
\qquad
\tau_0=10^4.
\]

\(h\) and \(\tau_0\) are not states specifying the physical case; they
are treated as constants of the numerical resolution and of the readout
encoding.

\subsection{6.3 The 41-real-component full
microstate}\label{the-41-real-component-full-microstate}

Of the eight persistent logical states, \(U,H\) are complex, so in the
implementation they become the 10 real components

\[
(U_R,U_I,P,E,H_R,H_I,Q,N,C,D).
\]

To these are added the RK4

\[
(k_1,k_2,k_3,k_4),
\]

the three components of each \(k_j\), the stage states, the composite
increment, the midpoint, \(\Delta\phi\) and the 11-phase one-hot phase

\[
q=(q_0,\ldots,q_{10}),
\]

making the full state 41 real components.

Conceptually,

\[
\begin{aligned}
z=&(U_R,U_I,P,E,H_R,H_I,Q,N,C,D;\\
&k_1,k_2,k_3,k_4;P_s,E_s;d;P_m,E_m;\Delta\phi;q_0,\ldots,q_{10}).
\end{aligned}
\]

Here the work states are not pushed outside on the grounds that ``they
are used only temporarily until the next macro state, so they are not
states''; if they are needed between microsteps they are made explicit
as part of the full state.

\subsection{6.4 Correspondence between the initialization code and the
formulas}\label{correspondence-between-the-initialization-code-and-the-formulas}

The initialization function actually used in the generator for the
reference case is the following. It is the actual code of
\texttt{run\_strict\_charged8\_final.py} itself.

\begin{Shaded}
\begin{Highlighting}[]
\KeywordTok{def}\NormalTok{ init\_state(p0}\OperatorTok{=}\FloatTok{50.0}\NormalTok{,n0}\OperatorTok{=}\FloatTok{0.25}\NormalTok{,c0}\OperatorTok{=}\FloatTok{1.09}\NormalTok{,d0}\OperatorTok{=}\FloatTok{0.36}\NormalTok{):}
\NormalTok{    z}\OperatorTok{=}\NormalTok{np.zeros(NST,dtype}\OperatorTok{=}\NormalTok{np.float64)}
\NormalTok{    z[UR]}\OperatorTok{=}\FloatTok{1.0}\OperatorTok{;}\NormalTok{ z[P]}\OperatorTok{=}\NormalTok{p0}\OperatorTok{;}\NormalTok{ z[HR]}\OperatorTok{=}\FloatTok{1.0}\OperatorTok{;}\NormalTok{ z[Q]}\OperatorTok{=}\FloatTok{1.0}
\NormalTok{    z[N]}\OperatorTok{=}\NormalTok{n0}\OperatorTok{;}\NormalTok{ z[C]}\OperatorTok{=}\NormalTok{c0}\OperatorTok{;}\NormalTok{ z[D]}\OperatorTok{=}\NormalTok{d0}
\NormalTok{    z[PS]}\OperatorTok{=}\NormalTok{p0}\OperatorTok{;}\NormalTok{ z[PM]}\OperatorTok{=}\NormalTok{p0}\OperatorTok{;}\NormalTok{ z[QB]}\OperatorTok{=}\FloatTok{1.0}
    \ControlFlowTok{return}\NormalTok{ z}
\end{Highlighting}
\end{Shaded}

By \texttt{np.zeros}, all components not explicitly assigned start from
0. For the persistent logical states this corresponds to

\[
U_0=1+0i,
\qquad
P_0=p_0,
\qquad
E_0=0,
\]

\[
H_0=1+0i,
\qquad
Q_0=1,
\]

\[
N_0=n_0,
\qquad
C_0=c_0,
\qquad
D_0=d_0.
\]

In the reference case

\[
p_0=50,
\qquad
n_0=0.25,
\qquad
c_0=1.09,
\qquad
d_0=0.36.
\]

Since

\[
Q_0=1=e^{0/\tau_0},
\]

the initial readout time is \(t_0=0\). Also,

\[
U_0=e^{i\chi_0}=1,
\qquad
H_0=e^{i\phi_0}=1
\]

express

\[
\chi_0=0,
\qquad
\phi_0=0.
\]

\texttt{z{[}PS{]}=p0} and \texttt{z{[}PM{]}=p0} are the initial values
of the work registers for the RK4 stage / midpoint, and
\texttt{z{[}QB{]}=1.0} expresses that the 11-phase one-hot phase is

\[
q=(1,0,\ldots,0),
\]

that is, the processing starts from phase 0.

What matters is that the individual

\[
\lambda_A,
\qquad
\lambda_B
\]

are not passed to \texttt{transition(z)}. The charge condition is
mapped, before the initialization, to

\[
C=1-\lambda_A\lambda_B,
\qquad
D=(\lambda_A-\lambda_B)^2,
\]

and thereafter \(C,D\) propagate as persistent states. Therefore, the
only physical information changed in the five-condition test is the
initial values of \(C,D\) in the same state format.

\subsection{6.5 Correspondence between the local interaction-rate code
and the
formulas}\label{correspondence-between-the-local-interaction-rate-code-and-the-formulas}

The local rate, evaluated uniquely from the present state, is collected
in the following function in the actual code.

\begin{Shaded}
\begin{Highlighting}[]
\AttributeTok{@njit}\NormalTok{(cache}\OperatorTok{=}\VariableTok{True}\NormalTok{)}
\KeywordTok{def}\NormalTok{ \_rates(p,n,c,d):}
\NormalTok{    rp}\OperatorTok{=}\NormalTok{math.sqrt(p)}\OperatorTok{;}\NormalTok{ rc}\OperatorTok{=}\NormalTok{math.sqrt(c)}
\NormalTok{    dp}\OperatorTok{={-}}\NormalTok{(}\FloatTok{4.0}\OperatorTok{/}\FloatTok{3.0}\NormalTok{)}\OperatorTok{*}\NormalTok{n}\OperatorTok{*}\NormalTok{d}\OperatorTok{*}\NormalTok{rc}\OperatorTok{/}\NormalTok{rp }\OperatorTok{{-}}\NormalTok{(}\FloatTok{64.0}\OperatorTok{/}\FloatTok{5.0}\NormalTok{)}\OperatorTok{*}\NormalTok{n}\OperatorTok{*}\NormalTok{c}\OperatorTok{*}\NormalTok{rc}\OperatorTok{/}\NormalTok{(p}\OperatorTok{*}\NormalTok{rp)}
\NormalTok{    dt}\OperatorTok{=}\NormalTok{p}\OperatorTok{*}\NormalTok{rp}\OperatorTok{/}\NormalTok{rc}
    \ControlFlowTok{return}\NormalTok{ dp,}\FloatTok{0.0}\NormalTok{,dt}
\end{Highlighting}
\end{Shaded}

These three components correspond directly to

\[
R(P,E,N,C,D)
=
\left(
\frac{dP}{d\phi},
\frac{dE}{d\phi},
\frac{dt}{d\phi}
\right).
\]

The line

\begin{Shaded}
\begin{Highlighting}[]
\NormalTok{dp}\OperatorTok{={-}}\NormalTok{(}\FloatTok{4.0}\OperatorTok{/}\FloatTok{3.0}\NormalTok{)}\OperatorTok{*}\NormalTok{n}\OperatorTok{*}\NormalTok{d}\OperatorTok{*}\NormalTok{rc}\OperatorTok{/}\NormalTok{rp }\OperatorTok{{-}}\NormalTok{(}\FloatTok{64.0}\OperatorTok{/}\FloatTok{5.0}\NormalTok{)}\OperatorTok{*}\NormalTok{n}\OperatorTok{*}\NormalTok{c}\OperatorTok{*}\NormalTok{rc}\OperatorTok{/}\NormalTok{(p}\OperatorTok{*}\NormalTok{rp)}
\end{Highlighting}
\end{Shaded}

evaluates

\[
\boxed{
\frac{dP}{d\phi}
=
-\frac{4}{3}ND\frac{\sqrt C}{\sqrt P}
-\frac{64}{5}NC\frac{\sqrt C}{P\sqrt P}
}
\]

as it is. The first term is the \(r^{-2}\)-type radial decay originating
in the leading electric-dipole radiation, and the second is the
\(r^{-3}\)-type radial decay originating in the leading gravitational
quadrupole radiation, converted into phase derivatives by
\(d\phi/dt=\sqrt{C/P^3}\).

The second component of \texttt{return\ dp,0.0,dt} is 0 because this
paper is restricted to circular orbits,

\[
E=0,
\qquad
\frac{dE}{d\phi}=0.
\]

Also,

\begin{Shaded}
\begin{Highlighting}[]
\NormalTok{dt}\OperatorTok{=}\NormalTok{p}\OperatorTok{*}\NormalTok{rp}\OperatorTok{/}\NormalTok{rc}
\end{Highlighting}
\end{Shaded}

corresponds to

\[
\boxed{
\frac{dt}{d\phi}
=
\frac{P\sqrt P}{\sqrt C}.
}
\]

This \texttt{\_rates} does not switch between the attractive and the
repulsive case, or between the GW and the EM term, by an external
branch. The difference goes into the state values \(C,D\) and is
evaluated continuously within the same formula. In particular, for
equal-sign equal charge-to-mass ratios

\[
D=0,
\]

so that only the EM dipole term becomes 0 automatically, without
rewriting the code.

\subsection{\texorpdfstring{6.6 The actual code of the single map
\texttt{transition(z)}}{6.6 The actual code of the single map transition(z)}}\label{the-actual-code-of-the-single-map-transitionz}

The one-hot sum of products of the candidate values is the following.

\begin{Shaded}
\begin{Highlighting}[]
\AttributeTok{@njit}\NormalTok{(cache}\OperatorTok{=}\VariableTok{True}\NormalTok{)}
\KeywordTok{def}\NormalTok{ \_blend(q,c):}
\NormalTok{    x}\OperatorTok{=}\FloatTok{0.0}
    \ControlFlowTok{for}\NormalTok{ j }\KeywordTok{in} \BuiltInTok{range}\NormalTok{(}\DecValTok{11}\NormalTok{): x }\OperatorTok{+=}\NormalTok{ q[j]}\OperatorTok{*}\NormalTok{c[j]}
    \ControlFlowTok{return}\NormalTok{ x}
\end{Highlighting}
\end{Shaded}

This implements the formula

\[
\boxed{
z_i'=\sum_{j=0}^{10}q_jF_{ij}(z)}
\]

as it is, as a sum of 11 products.

The core of the sole state update \texttt{transition(z)} used in the
reference case is the following.

\begin{Shaded}
\begin{Highlighting}[]
\AttributeTok{@njit}\NormalTok{(cache}\OperatorTok{=}\VariableTok{True}\NormalTok{)}
\KeywordTok{def}\NormalTok{ transition(z):}
    \CommentTok{"""Sole legal state{-}transition map."""}
\NormalTok{    q}\OperatorTok{=}\NormalTok{z[QB:QB}\OperatorTok{+}\DecValTok{11}\NormalTok{]}
\NormalTok{    p}\OperatorTok{=}\NormalTok{z[P]}\OperatorTok{;}\NormalTok{ e}\OperatorTok{=}\NormalTok{z[E]}\OperatorTok{;}\NormalTok{ n}\OperatorTok{=}\NormalTok{z[N]}\OperatorTok{;}\NormalTok{ c}\OperatorTok{=}\NormalTok{z[C]}\OperatorTok{;}\NormalTok{ dd}\OperatorTok{=}\NormalTok{z[D]}
    \CommentTok{\# old work values}
\NormalTok{    k1}\OperatorTok{=}\NormalTok{z[K1:K1}\OperatorTok{+}\DecValTok{3}\NormalTok{]}\OperatorTok{;}\NormalTok{ k2}\OperatorTok{=}\NormalTok{z[K2:K2}\OperatorTok{+}\DecValTok{3}\NormalTok{]}\OperatorTok{;}\NormalTok{ k3}\OperatorTok{=}\NormalTok{z[K3:K3}\OperatorTok{+}\DecValTok{3}\NormalTok{]}\OperatorTok{;}\NormalTok{ k4}\OperatorTok{=}\NormalTok{z[K4:K4}\OperatorTok{+}\DecValTok{3}\NormalTok{]}
\NormalTok{    ps}\OperatorTok{=}\NormalTok{z[PS]}\OperatorTok{;}\NormalTok{ es}\OperatorTok{=}\NormalTok{z[ES]}\OperatorTok{;}\NormalTok{ dv}\OperatorTok{=}\NormalTok{z[DV:DV}\OperatorTok{+}\DecValTok{3}\NormalTok{]}\OperatorTok{;}\NormalTok{ pm}\OperatorTok{=}\NormalTok{z[PM]}\OperatorTok{;}\NormalTok{ em}\OperatorTok{=}\NormalTok{z[EM]}\OperatorTok{;}\NormalTok{ dphi}\OperatorTok{=}\NormalTok{z[DPHI]}

    \CommentTok{\# stage rates are functions of OLD state only}
\NormalTok{    rP,rE,rT}\OperatorTok{=}\NormalTok{\_rates(p,n,c,dd)}
\NormalTok{    rPsP,rPsE,rPsT}\OperatorTok{=}\NormalTok{\_rates(ps,n,c,dd)}

\NormalTok{    cand}\OperatorTok{=}\NormalTok{np.empty((}\DecValTok{30}\NormalTok{,}\DecValTok{11}\NormalTok{),dtype}\OperatorTok{=}\NormalTok{np.float64)  }\CommentTok{\# all non{-}q registers 0..29}
    \CommentTok{\# default: identity candidate in all phases for every non{-}q register}
    \ControlFlowTok{for}\NormalTok{ i }\KeywordTok{in} \BuiltInTok{range}\NormalTok{(}\DecValTok{30}\NormalTok{):}
        \ControlFlowTok{for}\NormalTok{ j }\KeywordTok{in} \BuiltInTok{range}\NormalTok{(}\DecValTok{11}\NormalTok{): cand[i,j]}\OperatorTok{=}\NormalTok{z[i]}

    \CommentTok{\# k1: phase 0 compute, phase 10 reset}
\NormalTok{    cand[K1,}\DecValTok{0}\NormalTok{]}\OperatorTok{=}\NormalTok{rP}\OperatorTok{;}\NormalTok{ cand[K1}\OperatorTok{+}\DecValTok{1}\NormalTok{,}\DecValTok{0}\NormalTok{]}\OperatorTok{=}\NormalTok{rE}\OperatorTok{;}\NormalTok{ cand[K1}\OperatorTok{+}\DecValTok{2}\NormalTok{,}\DecValTok{0}\NormalTok{]}\OperatorTok{=}\NormalTok{rT}
    \ControlFlowTok{for}\NormalTok{ a }\KeywordTok{in} \BuiltInTok{range}\NormalTok{(}\DecValTok{3}\NormalTok{): cand[K1}\OperatorTok{+}\NormalTok{a,}\DecValTok{10}\NormalTok{]}\OperatorTok{=}\FloatTok{0.0}
    \CommentTok{\# stage2 fields phase 1}
\NormalTok{    cand[PS,}\DecValTok{1}\NormalTok{]}\OperatorTok{=}\NormalTok{p}\OperatorTok{+}\FloatTok{0.5}\OperatorTok{*}\NormalTok{HSTEP}\OperatorTok{*}\NormalTok{k1[}\DecValTok{0}\NormalTok{]}\OperatorTok{;}\NormalTok{ cand[ES,}\DecValTok{1}\NormalTok{]}\OperatorTok{=}\NormalTok{e}\OperatorTok{+}\FloatTok{0.5}\OperatorTok{*}\NormalTok{HSTEP}\OperatorTok{*}\NormalTok{k1[}\DecValTok{1}\NormalTok{]}
    \CommentTok{\# k2 phase 2, reset phase 10}
\NormalTok{    cand[K2,}\DecValTok{2}\NormalTok{]}\OperatorTok{=}\NormalTok{rPsP}\OperatorTok{;}\NormalTok{ cand[K2}\OperatorTok{+}\DecValTok{1}\NormalTok{,}\DecValTok{2}\NormalTok{]}\OperatorTok{=}\NormalTok{rPsE}\OperatorTok{;}\NormalTok{ cand[K2}\OperatorTok{+}\DecValTok{2}\NormalTok{,}\DecValTok{2}\NormalTok{]}\OperatorTok{=}\NormalTok{rPsT}
    \ControlFlowTok{for}\NormalTok{ a }\KeywordTok{in} \BuiltInTok{range}\NormalTok{(}\DecValTok{3}\NormalTok{): cand[K2}\OperatorTok{+}\NormalTok{a,}\DecValTok{10}\NormalTok{]}\OperatorTok{=}\FloatTok{0.0}
    \CommentTok{\# stage3 fields phase 3}
\NormalTok{    cand[PS,}\DecValTok{3}\NormalTok{]}\OperatorTok{=}\NormalTok{p}\OperatorTok{+}\FloatTok{0.5}\OperatorTok{*}\NormalTok{HSTEP}\OperatorTok{*}\NormalTok{k2[}\DecValTok{0}\NormalTok{]}\OperatorTok{;}\NormalTok{ cand[ES,}\DecValTok{3}\NormalTok{]}\OperatorTok{=}\NormalTok{e}\OperatorTok{+}\FloatTok{0.5}\OperatorTok{*}\NormalTok{HSTEP}\OperatorTok{*}\NormalTok{k2[}\DecValTok{1}\NormalTok{]}
    \CommentTok{\# k3 phase 4 uses OLD Ps,Es}
\NormalTok{    cand[K3,}\DecValTok{4}\NormalTok{]}\OperatorTok{=}\NormalTok{rPsP}\OperatorTok{;}\NormalTok{ cand[K3}\OperatorTok{+}\DecValTok{1}\NormalTok{,}\DecValTok{4}\NormalTok{]}\OperatorTok{=}\NormalTok{rPsE}\OperatorTok{;}\NormalTok{ cand[K3}\OperatorTok{+}\DecValTok{2}\NormalTok{,}\DecValTok{4}\NormalTok{]}\OperatorTok{=}\NormalTok{rPsT}
    \ControlFlowTok{for}\NormalTok{ a }\KeywordTok{in} \BuiltInTok{range}\NormalTok{(}\DecValTok{3}\NormalTok{): cand[K3}\OperatorTok{+}\NormalTok{a,}\DecValTok{10}\NormalTok{]}\OperatorTok{=}\FloatTok{0.0}
    \CommentTok{\# stage4 fields phase 5}
\NormalTok{    cand[PS,}\DecValTok{5}\NormalTok{]}\OperatorTok{=}\NormalTok{p}\OperatorTok{+}\NormalTok{HSTEP}\OperatorTok{*}\NormalTok{k3[}\DecValTok{0}\NormalTok{]}\OperatorTok{;}\NormalTok{ cand[ES,}\DecValTok{5}\NormalTok{]}\OperatorTok{=}\NormalTok{e}\OperatorTok{+}\NormalTok{HSTEP}\OperatorTok{*}\NormalTok{k3[}\DecValTok{1}\NormalTok{]}
    \CommentTok{\# k4 phase 6 uses OLD Ps,Es}
\NormalTok{    cand[K4,}\DecValTok{6}\NormalTok{]}\OperatorTok{=}\NormalTok{rPsP}\OperatorTok{;}\NormalTok{ cand[K4}\OperatorTok{+}\DecValTok{1}\NormalTok{,}\DecValTok{6}\NormalTok{]}\OperatorTok{=}\NormalTok{rPsE}\OperatorTok{;}\NormalTok{ cand[K4}\OperatorTok{+}\DecValTok{2}\NormalTok{,}\DecValTok{6}\NormalTok{]}\OperatorTok{=}\NormalTok{rPsT}
    \ControlFlowTok{for}\NormalTok{ a }\KeywordTok{in} \BuiltInTok{range}\NormalTok{(}\DecValTok{3}\NormalTok{): cand[K4}\OperatorTok{+}\NormalTok{a,}\DecValTok{10}\NormalTok{]}\OperatorTok{=}\FloatTok{0.0}
    \CommentTok{\# combined d phase 7; reset 10}
    \ControlFlowTok{for}\NormalTok{ a }\KeywordTok{in} \BuiltInTok{range}\NormalTok{(}\DecValTok{3}\NormalTok{):}
\NormalTok{        cand[DV}\OperatorTok{+}\NormalTok{a,}\DecValTok{7}\NormalTok{]}\OperatorTok{=}\NormalTok{(HSTEP}\OperatorTok{/}\FloatTok{6.0}\NormalTok{)}\OperatorTok{*}\NormalTok{(k1[a]}\OperatorTok{+}\FloatTok{2.0}\OperatorTok{*}\NormalTok{k2[a]}\OperatorTok{+}\FloatTok{2.0}\OperatorTok{*}\NormalTok{k3[a]}\OperatorTok{+}\NormalTok{k4[a])}
\NormalTok{        cand[DV}\OperatorTok{+}\NormalTok{a,}\DecValTok{10}\NormalTok{]}\OperatorTok{=}\FloatTok{0.0}
    \CommentTok{\# midpoint phase 8; commit{-}safe values phase 10}
\NormalTok{    cand[PM,}\DecValTok{8}\NormalTok{]}\OperatorTok{=}\NormalTok{p}\OperatorTok{+}\FloatTok{0.5}\OperatorTok{*}\NormalTok{dv[}\DecValTok{0}\NormalTok{]}\OperatorTok{;}\NormalTok{ cand[EM,}\DecValTok{8}\NormalTok{]}\OperatorTok{=}\NormalTok{e}\OperatorTok{+}\FloatTok{0.5}\OperatorTok{*}\NormalTok{dv[}\DecValTok{1}\NormalTok{]}
\NormalTok{    cand[PM,}\DecValTok{10}\NormalTok{]}\OperatorTok{=}\NormalTok{p}\OperatorTok{+}\NormalTok{dv[}\DecValTok{0}\NormalTok{]}\OperatorTok{;}\NormalTok{ cand[EM,}\DecValTok{10}\NormalTok{]}\OperatorTok{=}\NormalTok{e}\OperatorTok{+}\NormalTok{dv[}\DecValTok{1}\NormalTok{]}
    \CommentTok{\# dphi phase 9; reset 10. Circular model dphi=h exactly.}
\NormalTok{    cand[DPHI,}\DecValTok{9}\NormalTok{]}\OperatorTok{=}\NormalTok{HSTEP}\OperatorTok{;}\NormalTok{ cand[DPHI,}\DecValTok{10}\NormalTok{]}\OperatorTok{=}\FloatTok{0.0}
    \CommentTok{\# commit phase 10 for persistent states}
\NormalTok{    ur}\OperatorTok{=}\NormalTok{z[UR]}\OperatorTok{;}\NormalTok{ ui}\OperatorTok{=}\NormalTok{z[UI]}\OperatorTok{;}\NormalTok{ hr}\OperatorTok{=}\NormalTok{z[HR]}\OperatorTok{;}\NormalTok{ hi}\OperatorTok{=}\NormalTok{z[HI]}
\NormalTok{    cand[UR,}\DecValTok{10}\NormalTok{]}\OperatorTok{=}\NormalTok{ur}\OperatorTok{*}\NormalTok{PH\_RE}\OperatorTok{{-}}\NormalTok{ui}\OperatorTok{*}\NormalTok{PH\_IM}\OperatorTok{;}\NormalTok{ cand[UI,}\DecValTok{10}\NormalTok{]}\OperatorTok{=}\NormalTok{ur}\OperatorTok{*}\NormalTok{PH\_IM}\OperatorTok{+}\NormalTok{ui}\OperatorTok{*}\NormalTok{PH\_RE}
\NormalTok{    cand[P,}\DecValTok{10}\NormalTok{]}\OperatorTok{=}\NormalTok{p}\OperatorTok{+}\NormalTok{dv[}\DecValTok{0}\NormalTok{]}\OperatorTok{;}\NormalTok{ cand[E,}\DecValTok{10}\NormalTok{]}\OperatorTok{=}\NormalTok{e}\OperatorTok{+}\NormalTok{dv[}\DecValTok{1}\NormalTok{]}
\NormalTok{    cd}\OperatorTok{=}\NormalTok{math.cos(dphi)}\OperatorTok{;}\NormalTok{ sd}\OperatorTok{=}\NormalTok{math.sin(dphi)}
\NormalTok{    cand[HR,}\DecValTok{10}\NormalTok{]}\OperatorTok{=}\NormalTok{hr}\OperatorTok{*}\NormalTok{cd}\OperatorTok{{-}}\NormalTok{hi}\OperatorTok{*}\NormalTok{sd}\OperatorTok{;}\NormalTok{ cand[HI,}\DecValTok{10}\NormalTok{]}\OperatorTok{=}\NormalTok{hr}\OperatorTok{*}\NormalTok{sd}\OperatorTok{+}\NormalTok{hi}\OperatorTok{*}\NormalTok{cd}
\NormalTok{    cand[Q,}\DecValTok{10}\NormalTok{]}\OperatorTok{=}\NormalTok{z[Q]}\OperatorTok{*}\NormalTok{math.exp(dv[}\DecValTok{2}\NormalTok{]}\OperatorTok{/}\NormalTok{TAU0)}
    \CommentTok{\# N,C,D remain identity candidates in all phases via default rows.}
\NormalTok{    cand[PS,}\DecValTok{10}\NormalTok{]}\OperatorTok{=}\NormalTok{p}\OperatorTok{+}\NormalTok{dv[}\DecValTok{0}\NormalTok{]}\OperatorTok{;}\NormalTok{ cand[ES,}\DecValTok{10}\NormalTok{]}\OperatorTok{=}\NormalTok{e}\OperatorTok{+}\NormalTok{dv[}\DecValTok{1}\NormalTok{]}

\NormalTok{    out}\OperatorTok{=}\NormalTok{np.empty(NST,dtype}\OperatorTok{=}\NormalTok{np.float64)}
    \ControlFlowTok{for}\NormalTok{ i }\KeywordTok{in} \BuiltInTok{range}\NormalTok{(}\DecValTok{30}\NormalTok{): out[i]}\OperatorTok{=}\NormalTok{\_blend(q,cand[i])}
    \CommentTok{\# fixed cyclic linear map q\textquotesingle{} = C\_11 q (not external loop counter)}
\NormalTok{    out[QB]}\OperatorTok{=}\NormalTok{q[}\DecValTok{10}\NormalTok{]}
    \ControlFlowTok{for}\NormalTok{ j }\KeywordTok{in} \BuiltInTok{range}\NormalTok{(}\DecValTok{1}\NormalTok{,}\DecValTok{11}\NormalTok{): out[QB}\OperatorTok{+}\NormalTok{j]}\OperatorTok{=}\NormalTok{q[j}\OperatorTok{{-}}\DecValTok{1}\NormalTok{]}
    \ControlFlowTok{return}\NormalTok{ out}
\end{Highlighting}
\end{Shaded}

The meaning of each phase is

\[
\begin{array}{c|l}
0 & \text{evaluation of }k_1\\
1 & \text{stage state for }k_2\\
2 & \text{evaluation of }k_2\\
3 & \text{stage state for }k_3\\
4 & \text{evaluation of }k_3\\
5 & \text{stage state for }k_4\\
6 & \text{evaluation of }k_4\\
7 & \Delta=(h/6)(k_1+2k_2+2k_3+k_4)\\
8 & \text{midpoint work state}\\
9 & \Delta\phi=h\\
10 & \text{commit to the persistent states}
\end{array}
\]

For example, in phase 1,

\begin{Shaded}
\begin{Highlighting}[]
\NormalTok{cand[PS,}\DecValTok{1}\NormalTok{]}\OperatorTok{=}\NormalTok{p}\OperatorTok{+}\FloatTok{0.5}\OperatorTok{*}\NormalTok{HSTEP}\OperatorTok{*}\NormalTok{k1[}\DecValTok{0}\NormalTok{]}
\end{Highlighting}
\end{Shaded}

corresponds to

\[
P_s=P+\frac{h}{2}k_{1P},
\]

and in phase 7,

\begin{Shaded}
\begin{Highlighting}[]
\NormalTok{cand[DV}\OperatorTok{+}\NormalTok{a,}\DecValTok{7}\NormalTok{]}\OperatorTok{=}\NormalTok{(HSTEP}\OperatorTok{/}\FloatTok{6.0}\NormalTok{)}\OperatorTok{*}\NormalTok{(k1[a]}\OperatorTok{+}\FloatTok{2.0}\OperatorTok{*}\NormalTok{k2[a]}\OperatorTok{+}\FloatTok{2.0}\OperatorTok{*}\NormalTok{k3[a]}\OperatorTok{+}\NormalTok{k4[a])}
\end{Highlighting}
\end{Shaded}

corresponds to the usual RK4 composition

\[
\Delta
=
\frac{h}{6}
\left(k_1+2k_2+2k_3+k_4\right).
\]

Unlike ordinary sequential code, however, a next value written just
before within the same microstep is not used for the next operation. The
candidates of each phase are built from the OLD full state read at the
beginning of \texttt{transition(z)}, and only the candidate valid in
that microstep is selected by the one-hot \(q\). The next evaluation is
performed only after moving to the next microstep.

In the commit phase,

\begin{Shaded}
\begin{Highlighting}[]
\NormalTok{cand[P,}\DecValTok{10}\NormalTok{]}\OperatorTok{=}\NormalTok{p}\OperatorTok{+}\NormalTok{dv[}\DecValTok{0}\NormalTok{]}
\NormalTok{cand[E,}\DecValTok{10}\NormalTok{]}\OperatorTok{=}\NormalTok{e}\OperatorTok{+}\NormalTok{dv[}\DecValTok{1}\NormalTok{]}
\end{Highlighting}
\end{Shaded}

correspond to

\[
P'=P+\Delta P,
\qquad
E'=E+\Delta E.
\]

Also,

\begin{Shaded}
\begin{Highlighting}[]
\NormalTok{cand[Q,}\DecValTok{10}\NormalTok{]}\OperatorTok{=}\NormalTok{z[Q]}\OperatorTok{*}\NormalTok{math.exp(dv[}\DecValTok{2}\NormalTok{]}\OperatorTok{/}\NormalTok{TAU0)}
\end{Highlighting}
\end{Shaded}

is

\[
Q'=Q\exp\left(\frac{\Delta t}{\tau_0}\right),
\]

which maintains the clock encoding

\[
Q=e^{t/\tau_0}
\]

without an additive time update.

For circular orbits the phase increment of one macro step is

\[
\Delta\phi=h,
\]

so the update of \(H=e^{i\phi}\) is

\[
H'=He^{ih}.
\]

In the code the two real components \texttt{HR,HI} are updated as a
rotation matrix.

One macro step applies the sole \texttt{transition} exactly 11 times, as
in the following actual code.

\begin{Shaded}
\begin{Highlighting}[]
\AttributeTok{@njit}\NormalTok{(cache}\OperatorTok{=}\VariableTok{True}\NormalTok{)}
\KeywordTok{def}\NormalTok{ macrostep(z):}
    \CommentTok{\# exactly 11 applications of the sole transition, no alternate formula}
    \ControlFlowTok{for}\NormalTok{ \_ }\KeywordTok{in} \BuiltInTok{range}\NormalTok{(}\DecValTok{11}\NormalTok{): z}\OperatorTok{=}\NormalTok{transition(z)}
    \ControlFlowTok{return}\NormalTok{ z}
\end{Highlighting}
\end{Shaded}

The macrostep itself is therefore not a separate physical update rule,
but only the repetition that runs through the 11 phases once.

\subsection{\texorpdfstring{6.7 The conserved relational quantities
\(N,C,D\)}{6.7 The conserved relational quantities N,C,D}}\label{the-conserved-relational-quantities-ncd}

\texttt{cand} is first filled, for every non-\(q\) register, with the
identity candidate of the OLD state.

\begin{Shaded}
\begin{Highlighting}[]
\ControlFlowTok{for}\NormalTok{ i }\KeywordTok{in} \BuiltInTok{range}\NormalTok{(}\DecValTok{30}\NormalTok{):}
    \ControlFlowTok{for}\NormalTok{ j }\KeywordTok{in} \BuiltInTok{range}\NormalTok{(}\DecValTok{11}\NormalTok{):}
\NormalTok{        cand[i,j]}\OperatorTok{=}\NormalTok{z[i]}
\end{Highlighting}
\end{Shaded}

Into the rows of \(N,C,D\), no other candidate is written afterwards in
any phase. Therefore, even through the one-hot sum of products,

\[
\boxed{
N'=N,
\qquad
C'=C,
\qquad
D'=D.
}
\]

This is not a copying of Python external constants at each step.
\(N,C,D\) are also contained in the full state, and propagate as
identities as the result of passing through the same candidate rows and
the same one-hot selection as the other registers.

\subsection{6.8 Correspondence between the readout code and the
formulas}\label{correspondence-between-the-readout-code-and-the-formulas}

The actual code of the readout for observation is the following.

\begin{Shaded}
\begin{Highlighting}[]
\KeywordTok{def}\NormalTok{ readout(z):}
\NormalTok{    t}\OperatorTok{=}\NormalTok{TAU0}\OperatorTok{*}\NormalTok{math.log(z[Q])}
    \ControlFlowTok{return}\NormalTok{ t,z[P]}\OperatorTok{*}\NormalTok{z[HR],z[P]}\OperatorTok{*}\NormalTok{z[HI]}
\end{Highlighting}
\end{Shaded}

The first component is the inverse map

\[
\boxed{
t=\tau_0\ln Q}
\]

of

\[
Q=e^{t/\tau_0}.
\]

Since the present experiment is a quasi-circular orbit with

\[
E=0,
\]

the radius is

\[
r=P.
\]

Since \(H=e^{i\phi}=H_R+iH_I\),

\[
PH
=
P(\cos\phi+i\sin\phi)
=
x+iy,
\]

so that

\[
\boxed{
x=P H_R,
\qquad
y=P H_I}.
\]

The return value of \texttt{readout(z)} is therefore exactly

\[
(t,x,y).
\]

This function is not part of the state update rule. In the generation
loop it is called only when writing the log for storing the states, and
its return value is not passed to \texttt{transition(z)}.

Since the terminal \(P=20\) generally does not coincide with a macrostep
grid point, only after the generation ends is a linear interpolation
made from the two adjacent points,

\begin{Shaded}
\begin{Highlighting}[]
\NormalTok{t1,\_,\_}\OperatorTok{=}\NormalTok{readout(z)}\OperatorTok{;}\NormalTok{ p1}\OperatorTok{=}\NormalTok{z[P]}
\NormalTok{t0,\_,\_}\OperatorTok{=}\NormalTok{readout(prev)}\OperatorTok{;}\NormalTok{ p0}\OperatorTok{=}\NormalTok{prev[P]}
\NormalTok{frac}\OperatorTok{=}\NormalTok{(p0}\OperatorTok{{-}}\FloatTok{20.0}\NormalTok{)}\OperatorTok{/}\NormalTok{(p0}\OperatorTok{{-}}\NormalTok{p1)}
\NormalTok{tcross}\OperatorTok{=}\NormalTok{t0}\OperatorTok{+}\NormalTok{frac}\OperatorTok{*}\NormalTok{(t1}\OperatorTok{{-}}\NormalTok{t0)}
\NormalTok{phicross}\OperatorTok{=}\NormalTok{(nstep}\OperatorTok{{-}}\DecValTok{1}\OperatorTok{+}\NormalTok{frac)}\OperatorTok{*}\NormalTok{HSTEP}
\end{Highlighting}
\end{Shaded}

which corresponds to

\[
f
=
\frac{P_0-20}{P_0-P_1},
\]

\[
t_{20}=t_0+f(t_1-t_0),
\]

\[
\phi_{20}=(n-1+f)h.
\]

The \(t_{20},\phi_{20}\) obtained here are also readouts for comparison,
and are not fed back into the orbit already generated.

\subsection{6.9 Qualification concerning ``fully
multiplicative''}\label{qualification-concerning-fully-multiplicative}

The strict single map of this implementation means

\begin{itemize}
\tightlist
\item
  no external physical argument,
\item
  a sole \texttt{transition(z)},
\item
  OLD-state-only,
\item
  explicit work states,
\item
  one-hot sum of products,
\item
  non-feedback readout.
\end{itemize}

On the other hand, it is not claimed that square roots, reciprocals and
exponential functions have been removed from the primitive operations.
They remain as ``interaction functions computed uniquely from the
present state''. This paper is therefore not a proof that every
operation has been reduced to multiplication alone.

\begin{center}\rule{0.5\linewidth}{0.5pt}\end{center}

\section{7. Experimental design}\label{experimental-design}

\subsection{7.1 The reference case}\label{the-reference-case}

The reference case is

\[
\lambda_A=+0.30,
\qquad
\lambda_B=-0.30,
\]

\[
N=0.25,
\qquad
C=1.09,
\qquad
D=0.36,
\qquad
P_0=50.
\]

The initial state is

\[
U_0=1,
\qquad
H_0=1,
\qquad
Q_0=1,
\qquad
E_0=0.
\]

The generator continues the state update until \(P\le20\) is reached.

\subsection{7.2 The five-condition transferability
test}\label{the-five-condition-transferability-test}

From the series of integer multiples of the reference value \(0.30\),
the five cases satisfying \(C>0\) in the present quasi-circular
reference are chosen.

{\def\LTcaptype{none} % do not increment counter
\begin{longtable}[]{@{}
  >{\raggedright\arraybackslash}p{(\linewidth - 10\tabcolsep) * \real{0.1364}}
  >{\raggedleft\arraybackslash}p{(\linewidth - 10\tabcolsep) * \real{0.1818}}
  >{\raggedleft\arraybackslash}p{(\linewidth - 10\tabcolsep) * \real{0.1818}}
  >{\raggedleft\arraybackslash}p{(\linewidth - 10\tabcolsep) * \real{0.1818}}
  >{\raggedleft\arraybackslash}p{(\linewidth - 10\tabcolsep) * \real{0.1818}}
  >{\raggedright\arraybackslash}p{(\linewidth - 10\tabcolsep) * \real{0.1364}}@{}}
\toprule\noalign{}
\begin{minipage}[b]{\linewidth}\raggedright
case
\end{minipage} & \begin{minipage}[b]{\linewidth}\raggedleft
\(\lambda_A\)
\end{minipage} & \begin{minipage}[b]{\linewidth}\raggedleft
\(\lambda_B\)
\end{minipage} & \begin{minipage}[b]{\linewidth}\raggedleft
\(C\)
\end{minipage} & \begin{minipage}[b]{\linewidth}\raggedleft
\(D\)
\end{minipage} & \begin{minipage}[b]{\linewidth}\raggedright
Physical feature
\end{minipage} \\
\midrule\noalign{}
\endhead
\bottomrule\noalign{}
\endlastfoot
n1 attractive & +0.30 & -0.30 & 1.09 & 0.36 & Coulomb attraction + EM
dipole \\
n1 repulsive & +0.30 & +0.30 & 0.91 & 0 & Coulomb repulsion, no leading
dipole \\
n3 attractive & +0.90 & -0.90 & 1.81 & 3.24 & Strong Coulomb attraction
+ strong dipole \\
n3 repulsive & +0.90 & +0.90 & 0.19 & 0 & Strong Coulomb repulsion, no
leading dipole \\
n4 attractive & +1.20 & -1.20 & 2.44 & 5.76 & Strongest attractive side
+ dipole \\
\end{longtable}
}

For \(\lambda_A=\lambda_B=1.20\),

\[
C=1-1.44=-0.44,
\]

which is outside the domain of the present real quasi-circular analytic
reference, so it is not included in the five cases. Corrections such as
taking \(|C|\) would be a different physical model, and are not made.

What is changed in the five cases is only the \(C,D\) of the initial
state; \texttt{transition(z)}, \texttt{macrostep(z)}, the 11-phase
machine, the eight persistent states and the numerical resolution are
not changed.

\begin{center}\rule{0.5\linewidth}{0.5pt}\end{center}

\section{8. Results}\label{results}

\subsection{8.1 First, ``how the orbit changes with the
condition''}\label{first-how-the-orbit-changes-with-the-condition}

The most important visualization in Paper 6 is the orbit itself under
each condition.

\begin{figure}
\centering
\pandocbounded{\includegraphics[keepaspectratio,alt={Figure 1: Independent analytic reference orbits for the five conditions}]{figs/figure01_five_case_reference_orbits_same_scale.png}}
\caption{Figure 1: Independent analytic reference orbits for the five
conditions}
\end{figure}

\textbf{Figure 1.} The independent analytic reference orbits for the
five conditions. The display range of \(x,y\) is fixed to the same in
all panels. Each condition is drawn for \(r=50\rightarrow20\). For
opposite signs, the larger \(C,D\) the faster the inspiral; for
equal-sign equal ratios \(D=0\), so the leading EM dipole vanishes, and
in particular the n3 repulsive case with \(C=0.19\) requires a very
large number of orbits.

For the same five conditions, the overlays with the stored strict
eight-state generator are placed side by side.

\begin{figure}
\centering
\pandocbounded{\includegraphics[keepaspectratio,alt={Figure 2: Overlay of the analytic reference and the strict eight-state generator for the five conditions}]{figs/figure02_five_case_strict_generator_overlay_panels.png}}
\caption{Figure 2: Overlay of the analytic reference and the strict
eight-state generator for the five conditions}
\end{figure}

\textbf{Figure 2.} Orbit overlays of the analytic reference and the
strict 8-state generator for all five conditions. The source figure of
each panel is integrated as it is from the stored individual
experiments; the orbit data themselves are not altered.

Comparing the radius and the accumulated number of orbits directly, the
difference in density of the orbit figures can be read quantitatively.

\begin{figure}
\centering
\pandocbounded{\includegraphics[keepaspectratio,alt={Figure 3: Radius and accumulated number of orbits}]{figs/figure03_five_case_radius_vs_cycles.png}}
\caption{Figure 3: Radius and accumulated number of orbits}
\end{figure}

\textbf{Figure 3.} Radius and accumulated number of orbits for
\(r=50\rightarrow20\). The speed of the orbital decay changes by several
orders of magnitude over the five conditions.

Comparing only the first 10 orbits at the same spatial scale makes the
difference in the decay speed at the initial stage easy to read.

\begin{figure}
\centering
\pandocbounded{\includegraphics[keepaspectratio,alt={Figure 4: The first 10 orbits for the five conditions}]{figs/figure04_five_case_first_10_orbits_same_scale.png}}
\caption{Figure 4: The first 10 orbits for the five conditions}
\end{figure}

\textbf{Figure 4.} Comparison of the orbits over the first 10 orbits.
The display range is fixed, showing the difference of the initial
inspiral between the conditions.

The total number of orbits and the final time of the analytic reference
are as follows.

{\def\LTcaptype{none} % do not increment counter
\begin{longtable}[]{@{}
  >{\raggedright\arraybackslash}p{(\linewidth - 12\tabcolsep) * \real{0.1111}}
  >{\raggedleft\arraybackslash}p{(\linewidth - 12\tabcolsep) * \real{0.1481}}
  >{\raggedleft\arraybackslash}p{(\linewidth - 12\tabcolsep) * \real{0.1481}}
  >{\raggedleft\arraybackslash}p{(\linewidth - 12\tabcolsep) * \real{0.1481}}
  >{\raggedleft\arraybackslash}p{(\linewidth - 12\tabcolsep) * \real{0.1481}}
  >{\raggedleft\arraybackslash}p{(\linewidth - 12\tabcolsep) * \real{0.1481}}
  >{\raggedleft\arraybackslash}p{(\linewidth - 12\tabcolsep) * \real{0.1481}}@{}}
\toprule\noalign{}
\begin{minipage}[b]{\linewidth}\raggedright
case
\end{minipage} & \begin{minipage}[b]{\linewidth}\raggedleft
\(C\)
\end{minipage} & \begin{minipage}[b]{\linewidth}\raggedleft
\(D\)
\end{minipage} & \begin{minipage}[b]{\linewidth}\raggedleft
\(A_{\rm EM}\)
\end{minipage} & \begin{minipage}[b]{\linewidth}\raggedleft
\(B_{\rm GW}\)
\end{minipage} & \begin{minipage}[b]{\linewidth}\raggedleft
cycles \(50\to20\)
\end{minipage} & \begin{minipage}[b]{\linewidth}\raggedleft
\(t_f\)
\end{minipage} \\
\midrule\noalign{}
\endhead
\bottomrule\noalign{}
\endlastfoot
n1 attractive & 1.09 & 0.36 & 0.1308 & 3.80192 & 122.086006 &
169032.317607 \\
n1 repulsive & 0.91 & 0 & 0 & 2.64992 & 364.132605 & 574545.646661 \\
n3 attractive & 1.81 & 3.24 & 1.9548 & 10.48352 & 16.671907 &
17450.031850 \\
n3 repulsive & 0.19 & 0 & 0 & 0.11552 & 3816.728393 & 13179536.011080 \\
n4 attractive & 2.44 & 5.76 & 4.6848 & 19.05152 & 8.348877 &
7507.620157 \\
\end{longtable}
}

What this table shows is not the mere monotone relation ``stronger
charge, faster''. The direction of \(C\) changes with the sign, and for
equal-sign equal ratios \(D=0\), so that both the conservative force and
the dissipation change at the same time.

\subsection{8.2 Strict eight-state generation of the reference
case}\label{strict-eight-state-generation-of-the-reference-case}

The strict generation of the reference case reached \(P=20\) with

\begin{itemize}
\tightlist
\item
  macro steps: 488,345
\item
  microsteps: 5,371,795
\end{itemize}

The readouts of the generator are

\[
t_{\rm gen}=169032.317607049,
\]

\[
\phi_{\rm gen}=767.088999401507.
\]

The independent analytic reference is

\[
t_{\rm ref}=169032.31760706357,
\]

\[
\phi_{\rm ref}=767.0889994015913.
\]

Therefore

\[
\boxed{
|\Delta t|=1.463923\times10^{-8}
}
\]

\[
\boxed{
|\Delta\phi|=8.435563\times10^{-11}
}
\]

In the comparison at identical physical times over the whole stored
range,

\[
\max|\Delta r|=9.116260\times10^{-9},
\]

\[
\operatorname{RMS}(|\Delta r|)=3.771324\times10^{-9},
\]

\[
\max\|\Delta(x,y)\|=3.763773\times10^{-6},
\]

\[
\operatorname{RMS}(\|\Delta(x,y)\|)=2.130252\times10^{-6}.
\]

\begin{figure}
\centering
\pandocbounded{\includegraphics[keepaspectratio,alt={Figure 5: Radius-time comparison for the reference case}]{figs/figure05_baseline_radius_vs_time.png}}
\caption{Figure 5: Radius-time comparison for the reference case}
\end{figure}

\textbf{Figure 5.} Radius--time comparison for the reference case. The
independent analytic reference is used only for the comparison after the
state generation.

\begin{figure}
\centering
\pandocbounded{\includegraphics[keepaspectratio,alt={Figure 6: Residuals for the reference case}]{figs/figure06_baseline_reference_residuals.png}}
\caption{Figure 6: Residuals for the reference case}
\end{figure}

\textbf{Figure 6.} Residuals against the independent analytic reference
in the reference case.

\begin{figure}
\centering
\pandocbounded{\includegraphics[keepaspectratio,alt={Figure 7: Identity holding of N, C, D}]{figs/figure07_identity_states_NCD.png}}
\caption{Figure 7: Identity holding of N, C, D}
\end{figure}

\textbf{Figure 7.} Holding of \(N,C,D\) in the reference case. The
maximum drift over the whole run was 0.

\begin{figure}
\centering
\pandocbounded{\includegraphics[keepaspectratio,alt={Figure 8: The first macrostep of the 11-phase machine/work states}]{figs/figure08_microphase_work_states.png}}
\caption{Figure 8: The first macrostep of the 11-phase machine/work
states}
\end{figure}

\textbf{Figure 8.} Evolution of the explicit work states in the first
macrostep. The RK4 stage values are not reused directly within the same
microstep; they are referred to for the first time as OLD state in the
next microstep.

\begin{figure}
\centering
\pandocbounded{\includegraphics[keepaspectratio,alt={Figure 9: EM dipole / GW quadrupole power ratio in the reference case}]{figs/figure09_em_to_gw_power_ratio.png}}
\caption{Figure 9: EM dipole / GW quadrupole power ratio in the
reference case}
\end{figure}

\textbf{Figure 9.} \(P_{\rm EM}/P_{\rm GW}\) in the reference case. This
figure is a diagnostic on the side of the independent analytic reference
and is not fed back into the update of the state generator.

\subsection{8.3 All five conditions generated by the same strict
map}\label{all-five-conditions-generated-by-the-same-strict-map}

For the five conditions of Figure 2, the numerical summary is as
follows.

{\def\LTcaptype{none} % do not increment counter
\begin{longtable}[]{@{}
  >{\raggedright\arraybackslash}p{(\linewidth - 10\tabcolsep) * \real{0.1429}}
  >{\raggedleft\arraybackslash}p{(\linewidth - 10\tabcolsep) * \real{0.1905}}
  >{\raggedleft\arraybackslash}p{(\linewidth - 10\tabcolsep) * \real{0.1905}}
  >{\raggedleft\arraybackslash}p{(\linewidth - 10\tabcolsep) * \real{0.1905}}
  >{\raggedright\arraybackslash}p{(\linewidth - 10\tabcolsep) * \real{0.1429}}
  >{\raggedright\arraybackslash}p{(\linewidth - 10\tabcolsep) * \real{0.1429}}@{}}
\toprule\noalign{}
\begin{minipage}[b]{\linewidth}\raggedright
case
\end{minipage} & \begin{minipage}[b]{\linewidth}\raggedleft
macro steps
\end{minipage} & \begin{minipage}[b]{\linewidth}\raggedleft
\(|\Delta\phi|\)
\end{minipage} & \begin{minipage}[b]{\linewidth}\raggedleft
\(|\Delta t|\)
\end{minipage} & \begin{minipage}[b]{\linewidth}\raggedright
time readout
\end{minipage} & \begin{minipage}[b]{\linewidth}\raggedright
\(\Delta N,\Delta C,\Delta D\)
\end{minipage} \\
\midrule\noalign{}
\endhead
\bottomrule\noalign{}
\endlastfoot
n1 attractive & 488,345 & \(8.436\times10^{-11}\) &
\(1.464\times10^{-8}\) & finite & 0, 0, 0 \\
n1 repulsive & 1,456,531 & \(1.305\times10^{-9}\) &
\(3.203\times10^{-7}\) & finite & 0, 0, 0 \\
n3 attractive & 66,688 & \(4.229\times10^{-9}\) & \(8.739\times10^{-7}\)
& finite & 0, 0, 0 \\
n3 repulsive & 15,266,914 & \(3.653\times10^{-9}\) & N/A & nonfinite
after step 6,316,005 & 0, 0, 0 \\
n4 attractive & 33,396 & \(8.323\times10^{-9}\) & \(1.545\times10^{-6}\)
& finite & 0, 0, 0 \\
\end{longtable}
}

What matters is that no separate generating formula was used for the
five cases. In all cases,

\[
\boxed{
\text{same }transition(z)
+
\text{same 11-phase machine}
+
\text{same }Z_8
}
\]

were used, and the difference of the physical condition was given only
by the initial \(C,D\) of the state.

\subsection{8.4 The time-readout limit of n3
repulsive}\label{the-time-readout-limit-of-n3-repulsive}

\texttt{n3\_repulsive} has

\[
C=0.19,
\qquad
D=0,
\]

and is the slowest case, requiring about 3816.7 orbits.

The orbital states \(P,H\) could be generated to the end, and the phase
residual is finite. On the other hand, the clock encoding by

\[
Q=e^{t/\tau_0}
\]

exceeded the finite-precision range because of the long duration, and
became nonfinite after step 6,316,005.

In this paper this case is not deleted, and no clipping of \(Q\),
renormalization, or switch to another clock is made. This is recorded as

\begin{center}
\fbox{\parbox{0.90\linewidth}{\centering not a failure of the orbit generation, but a finite-precision limit of the present exponential clock readout$.$}}
\end{center}

However, since the time cannot be read out to the end, no agreement of a
finite value of \(\Delta t\) is claimed for this case.

\begin{center}\rule{0.5\linewidth}{0.5pt}\end{center}

\section{9. Discussion}\label{discussion}

\subsection{9.1 Not the individual charges but relational quantities
remained in the
state}\label{not-the-individual-charges-but-relational-quantities-remained-in-the-state}

In the present restricted model, what \texttt{transition(z)} directly
needs is not

\[
\lambda_A,
\lambda_B
\]

themselves, but

\[
C=1-\lambda_A\lambda_B,
\qquad
D=(\lambda_A-\lambda_B)^2.
\]

This shows that the state can be expressed by the \textbf{relational
quantities} that actually appear in the interaction formulas, rather
than by labelled parameters attached to individual bodies. This is,
however, a representational fact in the present quasi-circular LO model;
it does not prove that \(C,D\) alone are fundamental in general charged
binaries too.

\subsection{9.2 Attraction/repulsion was not made an external
branch}\label{attractionrepulsion-was-not-made-an-external-branch}

\texttt{transition(z)} contains

\begin{itemize}
\tightlist
\item
  no attractive / repulsive if-branch,
\item
  no identity of bodies A/B,
\item
  no gravity / Coulomb force-type branch.
\end{itemize}

The difference of the state values, \(C>1,D>0\) for opposite signs and
\(C<1,D=0\) for equal-sign equal ratios, unfolds as it is into different
orbits.

What was confirmed in this experiment is therefore

\begin{center}
\fbox{\parbox{0.90\linewidth}{\centering that different physical conditions can be handled not by switching algorithms but as differences of state values$.$}}
\end{center}

\subsection{9.3 Two kinds of dissipation enter one
rate}\label{two-kinds-of-dissipation-enter-one-rate}

The radius evolution is

\[
\frac{dr}{dt}
=-\frac{A}{r^2}-\frac{B}{r^3},
\]

and the EM dipole and the GW quadrupole have different distance
dependences. Nevertheless, converted into phase derivatives,

\[
\frac{dP}{d\phi}
=
-\frac{4}{3}ND\frac{\sqrt C}{\sqrt P}
-\frac{64}{5}NC\frac{\sqrt C}{P\sqrt P},
\]

they enter the same rate function, and \texttt{transition(z)} itself
need not be switched.

In this sense this paper gives \textbf{an example in which several known
interaction and dissipation terms can be internalized into one closed
state-update form}.

\subsection{9.4 This is not ``a unification of gravity and
electromagnetism''}\label{this-is-not-a-unification-of-gravity-and-electromagnetism}

That two kinds of contribution could be handled by the same state map,
and that they are the same fundamental interaction in nature, are
separate matters.

In this paper the known gravitational and electromagnetic terms are
given in advance as the analytic reference, and their coefficients and
the required states are audited. Therefore

\[
\boxed{
\text{same representation}
\not\Rightarrow
\text{same fundamental interaction}.
}
\]

\subsection{9.5 Eight states are not necessarily
minimal}\label{eight-states-are-not-necessarily-minimal}

Since \(C,D\) are constructed from \(\lambda_A,\lambda_B\), they are not
independent fundamental degrees of freedom. Also, in the present
symmetric series \(|\lambda_A|=|\lambda_B|\), so that the admissible
\((C,D)\) are strongly related.

Therefore it is not claimed that

\[
\boxed{
8\text{ is the mathematically minimal number of states}.
}
\]

What is confirmed now is only that

\[
\boxed{
\text{the present representation can be closed with }Z_8.
}
\]

\begin{center}\rule{0.5\linewidth}{0.5pt}\end{center}

\section{10. Limits and non-claims}\label{limits-and-non-claims}

This paper does not claim the following.

\begin{enumerate}
\def\labelenumi{\arabic{enumi}.}
\item
  \textbf{That the Einstein--Maxwell equations have been derived from
  eight states.}\\
  This paper reverse-constructs a known leading-order analytic model.
\item
  \textbf{That gravity and electromagnetism have been fundamentally
  unified.}\\
  That they can be handled in one state-update form is separate from the
  ontological identity of the interactions.
\item
  \textbf{That eight states are minimal or physically real.}\\
  \(Z_8\) is the set of persistent logical states that closes the
  present representation.
\item
  \textbf{Universality for general charged binaries.}\\
  Eccentricity, spin, tidal/finite-size effects, generic mass ratios,
  higher PN, NLO radiation, strong-field merger and full nonlinear
  Einstein--Maxwell numerical relativity are not treated.
\item
  \textbf{Quasi-circular orbits with \(C\le0\).}\\
  The present real quasi-circular reference requires \(C>0\).
\item
  \textbf{That all electromagnetic radiation vanishes when the leading
  dipole vanishes.}\\
  What \(D=0\) removes is the leading electric-dipole term adopted in
  this paper.
\item
  \textbf{That every internal operation has been reduced to
  multiplication alone.}\\
  Square roots, reciprocals, exponential functions and so on remain as
  current-state interaction functions.
\item
  \textbf{That the 11 phases are a fundamental structure of nature.}\\
  The 11 phases are the machine schedule that expands the present RK4
  processing into OLD-state-only form.
\end{enumerate}

Also, the nonfinite \(Q\) of \texttt{n3\_repulsive} is left as it is, as
a numerical limit of the present readout encoding.

\begin{center}\rule{0.5\linewidth}{0.5pt}\end{center}

\section{11. Reproducibility}\label{reproducibility}

(The folder names in the repository are in Japanese; they are
transliterated here.)

\subsection{11.1 Analytic reference}\label{analytic-reference}

The analytic formulas, the initial conditions, the independent solver,
the validation and the raw CSV files are stored in

\texttt{fifth\_\allowbreak{}thought\_\allowbreak{}experiment\_\allowbreak{}multiplicative\_\allowbreak{}state\_\allowbreak{}expansion\_\allowbreak{}and\_\allowbreak{}order\_\allowbreak{}verification\_\allowbreak{}20260924/\allowbreak{}07\_\allowbreak{}charged\_\allowbreak{}binary\_\allowbreak{}analytic\_\allowbreak{}reference\_\allowbreak{}20260925/\allowbreak{}}

Main code:

\begin{itemize}
\tightlist
\item
  \texttt{solve\_charged\_binary\_analytic\_reference\_v1.py}
\item
  \texttt{validate\_charged\_binary\_analytic\_reference\_v1.py}
\item
  \texttt{plot\_charged\_binary\_analytic\_reference\_v1.py}
\end{itemize}

The state generator does not read these reference data.

\subsection{11.2 Strict eight-state generator for the reference
case}\label{strict-eight-state-generator-for-the-reference-case}

Stored in:

\texttt{fifth\_\allowbreak{}thought\_\allowbreak{}experiment\_\allowbreak{}multiplicative\_\allowbreak{}state\_\allowbreak{}expansion\_\allowbreak{}and\_\allowbreak{}order\_\allowbreak{}verification\_\allowbreak{}20260924/\allowbreak{}11\_\allowbreak{}charged8\_\allowbreak{}strict\_\allowbreak{}single\_\allowbreak{}map\_\allowbreak{}experiment\_\allowbreak{}20260925/\allowbreak{}}

Main code:

\begin{itemize}
\tightlist
\item
  \texttt{01\_\allowbreak{}CODE/\allowbreak{}run\_\allowbreak{}strict\_\allowbreak{}charged8\_\allowbreak{}final.py}
\item
  \texttt{01\_\allowbreak{}CODE/\allowbreak{}postprocess\_\allowbreak{}strict\_\allowbreak{}charged8.py}
\end{itemize}

Audit evidence:

\begin{itemize}
\tightlist
\item
  \texttt{00\_\allowbreak{}RULES\_\allowbreak{}AND\_\allowbreak{}AUDIT/\allowbreak{}preexec\_\allowbreak{}audit.json}
\item
  \texttt{03\_\allowbreak{}ANALYSIS/\allowbreak{}generator\_\allowbreak{}summary.json}
\item
  \texttt{03\_\allowbreak{}ANALYSIS/\allowbreak{}reference\_\allowbreak{}comparison.json}
\end{itemize}

\subsection{11.3 Five-case transfer
experiment}\label{five-case-transfer-experiment}

Stored in:

\texttt{fifth\_\allowbreak{}thought\_\allowbreak{}experiment\_\allowbreak{}multiplicative\_\allowbreak{}state\_\allowbreak{}expansion\_\allowbreak{}and\_\allowbreak{}order\_\allowbreak{}verification\_\allowbreak{}20260924/\allowbreak{}13\_\allowbreak{}charged8\_\allowbreak{}integer\_\allowbreak{}multiple\_\allowbreak{}5case\_\allowbreak{}strict\_\allowbreak{}numeric\_\allowbreak{}20260925/\allowbreak{}}

Main code:

\begin{itemize}
\tightlist
\item
  \texttt{00\_\allowbreak{}CODE\_\allowbreak{}AND\_\allowbreak{}RULES/\allowbreak{}run\_\allowbreak{}strict\_\allowbreak{}charged8\_\allowbreak{}5cases.py}
\item
  \texttt{00\_\allowbreak{}CODE\_\allowbreak{}AND\_\allowbreak{}RULES/\allowbreak{}run\_\allowbreak{}n3\_\allowbreak{}repulsive\_\allowbreak{}checkpointed.py}
\item
  \texttt{00\_\allowbreak{}CODE\_\allowbreak{}AND\_\allowbreak{}RULES/\allowbreak{}postprocess\_\allowbreak{}strict\_\allowbreak{}charged8\_\allowbreak{}5cases.py}
\item
  \texttt{00\_\allowbreak{}CODE\_\allowbreak{}AND\_\allowbreak{}RULES/\allowbreak{}CASE\_\allowbreak{}MATRIX.json}
\end{itemize}

All macrostep rows are stored in HDF5 part files. \texttt{n3\_repulsive}
was run in parts because of the run-time limit, but what was carried
over between parts was only the complete 41-component full state; no
switch to another generating formula was made.

\subsection{11.4 Orbit comparison figures for Paper
6}\label{orbit-comparison-figures-for-paper-6}

Stored in:

\texttt{sixth\_\allowbreak{}thought\_\allowbreak{}experiment\_\allowbreak{}charged8\_\allowbreak{}gravity\_\allowbreak{}coulomb\_\allowbreak{}multiple\_\allowbreak{}interactions\_\allowbreak{}20260926/\allowbreak{}01\_\allowbreak{}orbit\_\allowbreak{}condition\_\allowbreak{}comparison\_\allowbreak{}figures/\allowbreak{}}

Reproduction code:

\begin{itemize}
\tightlist
\item
  \texttt{generate\_paper6\_orbit\_condition\_comparison\_v1.py}
\end{itemize}

Products:

\begin{itemize}
\tightlist
\item
  \texttt{figures/\allowbreak{}figure01\_\allowbreak{}five\_\allowbreak{}case\_\allowbreak{}reference\_\allowbreak{}orbits\_\allowbreak{}same\_\allowbreak{}scale.svg}
\item
  \texttt{figures/\allowbreak{}figure02\_\allowbreak{}five\_\allowbreak{}case\_\allowbreak{}strict\_\allowbreak{}generator\_\allowbreak{}overlay\_\allowbreak{}panels.svg}
\item
  \texttt{figures/\allowbreak{}figure03\_\allowbreak{}five\_\allowbreak{}case\_\allowbreak{}radius\_\allowbreak{}vs\_\allowbreak{}cycles.svg}
\item
  \texttt{figures/\allowbreak{}figure04\_\allowbreak{}five\_\allowbreak{}case\_\allowbreak{}first\_\allowbreak{}10\_\allowbreak{}orbits\_\allowbreak{}same\_\allowbreak{}scale.svg}
\item
  \texttt{figures/\allowbreak{}paper6\_\allowbreak{}orbit\_\allowbreak{}condition\_\allowbreak{}case\_\allowbreak{}summary.csv}
\end{itemize}

The source figures of the strict overlays of the five cases are stored
in \texttt{strict\_\allowbreak{}sources/\allowbreak{}}; the orbit data were not changed when
integrating them.

\subsection{11.5 Auxiliary figures for the main
text}\label{auxiliary-figures-for-the-main-text}

Stored in:

\texttt{sixth\_\allowbreak{}thought\_\allowbreak{}experiment\_\allowbreak{}charged8\_\allowbreak{}gravity\_\allowbreak{}coulomb\_\allowbreak{}multiple\_\allowbreak{}interactions\_\allowbreak{}20260926/\allowbreak{}02\_\allowbreak{}paper\_\allowbreak{}main\_\allowbreak{}figures/\allowbreak{}}

\begin{itemize}
\tightlist
\item
  \texttt{figure05\_baseline\_radius\_vs\_time.svg}
\item
  \texttt{figure06\_baseline\_reference\_residuals.svg}
\item
  \texttt{figure07\_identity\_states\_NCD.svg}
\item
  \texttt{figure08\_microphase\_work\_states.svg}
\item
  \texttt{figure09\_em\_to\_gw\_power\_ratio.svg}
\end{itemize}

\begin{center}\rule{0.5\linewidth}{0.5pt}\end{center}

\section{12. Conclusion}\label{conclusion}

In the fifth thought experiment {[}5{]}, the gravitational-orbit
generator was self-audited, and it was confirmed that, in the present
representation, not the apparent two states but six persistent logical
states and explicit machine/work states are needed.

In the sixth thought experiment, those audit rules were extended from
the gravity-only system to the charged binary with gravity + Coulomb.

Organizing the independent analytic reference, the charge dependence in
the present leading-order quasi-circular model was collected into

\[
C=1-\lambda_A\lambda_B
\]

and

\[
D=(\lambda_A-\lambda_B)^2.
\]

Adding these to the six states of the fifth thought experiment,

\[
\boxed{
Z_8=(U,P,E,H,Q,N,C,D),
}
\]

the charged quasi-circular inspiral could be generated without bringing
the external \(\lambda_A,\lambda_B\), an attraction/repulsion branch, or
a force-type branch into \texttt{transition(z)}.

In the reference case, the differences from the analytic reference were

\[
|\Delta t|\simeq1.46\times10^{-8},
\qquad
|\Delta\phi|\simeq8.44\times10^{-11},
\]

with

\[
\Delta N=\Delta C=\Delta D=0.
\]

Furthermore, with the same \texttt{transition(z)}, the same 11-phase
machine and the same eight persistent states, it could be transferred to
five kinds of charge condition. The difference of the conditions
unfolded into differences of orbit not as an external branch of the same
algorithm but as the values of \(C,D\) in the initial state.

The result obtained in this paper is therefore

\[
\boxed{
\begin{array}{c}
\text{for the restricted known gravity + Coulomb quasi-circular binary,}\\
\text{the conservative force and the two kinds of radiation reaction}\\
\text{can be internalized and reconstructed into a single synchronous state map}\\
\text{with no external physical parameters.}
\end{array}
}
\]

This is not ``the completion of a unified theory''. It is rather the
next fixed point of this research series, which \textbf{audits, one step
at a time, what information must be placed inside the state for the
system to close when known physics is reverse-constructed}.

\begin{center}\rule{0.5\linewidth}{0.5pt}\end{center}

\section{Appendix A. Audit items of the strict single
map}\label{appendix-a.-audit-items-of-the-strict-single-map}

In the pre-execution audit, the following were confirmed.

{\def\LTcaptype{none} % do not increment counter
\begin{longtable}[]{@{}
  >{\raggedright\arraybackslash}p{(\linewidth - 4\tabcolsep) * \real{0.3333}}
  >{\raggedright\arraybackslash}p{(\linewidth - 4\tabcolsep) * \real{0.3333}}
  >{\raggedright\arraybackslash}p{(\linewidth - 4\tabcolsep) * \real{0.3333}}@{}}
\toprule\noalign{}
\begin{minipage}[b]{\linewidth}\raggedright
Item
\end{minipage} & \begin{minipage}[b]{\linewidth}\raggedright
Content
\end{minipage} & \begin{minipage}[b]{\linewidth}\raggedright
Result
\end{minipage} \\
\midrule\noalign{}
\endhead
\bottomrule\noalign{}
\endlastfoot
A1 & The only argument of \texttt{transition} is \texttt{z}; the
persistent states are \(U,P,E,H,Q,N,C,D\) & pass \\
A2 & The sole generator is \texttt{transition}; no
fast/collapsed/shortcut path & pass \\
A3 & The non-\(q\) next state is the one-hot sum of products of 11
candidates; \(q\) is a fixed cyclic linear map & pass \\
A4 & Every candidate uses the OLD full state only; the work registers
are explicit & pass \\
A5 & The readout is at output only and is not fed back into the dynamics
& pass \\
A6 & No A/B identity, no force-type branch & pass \\
A7 & The physical case values \(N,C,D\) are states; the others are
resolution, encoding and RK coefficients & pass \\
A8 & The old failed experiments are stored; no reference import, no
other generator & pass \\
\end{longtable}
}

The failed experiments of the old \texttt{09\_...} and \texttt{10\_...}
are stored without deletion or overwriting.

\begin{center}\rule{0.5\linewidth}{0.5pt}\end{center}

\section{Appendix B. List of figures}\label{appendix-b.-list-of-figures}

\begin{enumerate}
\def\labelenumi{\arabic{enumi}.}
\item
  \textbf{Figure 1}
  \texttt{01\_\allowbreak{}orbit\_\allowbreak{}condition\_\allowbreak{}comparison\_\allowbreak{}figures/\allowbreak{}figures/\allowbreak{}figure01\_\allowbreak{}five\_\allowbreak{}case\_\allowbreak{}reference\_\allowbreak{}orbits\_\allowbreak{}same\_\allowbreak{}scale.svg}\\
  Independent analytic reference orbits for the five conditions, same
  scale.
\item
  \textbf{Figure 2}
  \texttt{01\_\allowbreak{}orbit\_\allowbreak{}condition\_\allowbreak{}comparison\_\allowbreak{}figures/\allowbreak{}figures/\allowbreak{}figure02\_\allowbreak{}five\_\allowbreak{}case\_\allowbreak{}strict\_\allowbreak{}generator\_\allowbreak{}overlay\_\allowbreak{}panels.svg}\\
  Analytic reference vs strict eight-state generator for the five
  conditions.
\item
  \textbf{Figure 3}
  \texttt{01\_\allowbreak{}orbit\_\allowbreak{}condition\_\allowbreak{}comparison\_\allowbreak{}figures/\allowbreak{}figures/\allowbreak{}figure03\_\allowbreak{}five\_\allowbreak{}case\_\allowbreak{}radius\_\allowbreak{}vs\_\allowbreak{}cycles.svg}\\
  Radius vs accumulated number of orbits.
\item
  \textbf{Figure 4}
  \texttt{01\_\allowbreak{}orbit\_\allowbreak{}condition\_\allowbreak{}comparison\_\allowbreak{}figures/\allowbreak{}figures/\allowbreak{}figure04\_\allowbreak{}five\_\allowbreak{}case\_\allowbreak{}first\_\allowbreak{}10\_\allowbreak{}orbits\_\allowbreak{}same\_\allowbreak{}scale.svg}\\
  Same-scale comparison of the first 10 orbits.
\item
  \textbf{Figure 5}
  \texttt{02\_\allowbreak{}paper\_\allowbreak{}main\_\allowbreak{}figures/\allowbreak{}figure05\_\allowbreak{}baseline\_\allowbreak{}radius\_\allowbreak{}vs\_\allowbreak{}time.svg}\\
  Radius--time comparison for the reference case.
\item
  \textbf{Figure 6}
  \texttt{02\_\allowbreak{}paper\_\allowbreak{}main\_\allowbreak{}figures/\allowbreak{}figure06\_\allowbreak{}baseline\_\allowbreak{}reference\_\allowbreak{}residuals.svg}\\
  Residuals against the independent analytic reference for the reference
  case.
\item
  \textbf{Figure 7}
  \texttt{02\_\allowbreak{}paper\_\allowbreak{}main\_\allowbreak{}figures/\allowbreak{}figure07\_\allowbreak{}identity\_\allowbreak{}states\_\allowbreak{}NCD.svg}\\
  Identity holding of \(N,C,D\).
\item
  \textbf{Figure 8}
  \texttt{02\_\allowbreak{}paper\_\allowbreak{}main\_\allowbreak{}figures/\allowbreak{}figure08\_\allowbreak{}microphase\_\allowbreak{}work\_\allowbreak{}states.svg}\\
  11-phase machine/work states.
\item
  \textbf{Figure 9}
  \texttt{02\_\allowbreak{}paper\_\allowbreak{}main\_\allowbreak{}figures/\allowbreak{}figure09\_\allowbreak{}em\_\allowbreak{}to\_\allowbreak{}gw\_\allowbreak{}power\_\allowbreak{}ratio.svg}\\
  EM/GW power ratio for the reference case.
\end{enumerate}

\begin{center}\rule{0.5\linewidth}{0.5pt}\end{center}

\section{References}\label{references}

\subsection{This series}\label{this-series}

{[}0{]} Noriaki Kihara, \textbf{Designing a Research Project to Explore
the Emergence of Spacetime, Particles, and Interactions from States and
Relations: A Framework of Preliminary Design and Feasibility
Verification for Foundational-Physics Model Exploration under
Uncertainty}, v1.0 (2026-09-20). Concept DOI: 10.5281/zenodo.22851944;
Version DOI: 10.5281/zenodo.22851945.

{[}1{]} Noriaki Kihara, \textbf{The First Thought Experiment: Tracing
the Equivalence Principle Down to Two Nameless Degrees of Freedom -
Deriving the Sign of the Sum of Squares, the Curvature Radius, and Three
Systems (Inertial, Uniformly Accelerated, Rotating) from the
Conservation of an Area Readout}, v1.0 (2026-09-20). Concept DOI:
10.5281/zenodo.22857952; Version DOI: 10.5281/zenodo.22857953.

{[}2{]} Noriaki Kihara, \textbf{The Second Thought Experiment: Deriving
the Coulomb-Type Inverse-Square Force from a Product Readout and an Area
Clock - Keeping Two Values and a Linear Law, Obtaining the Focus,
Attraction and Repulsion and Neutrality, and Fixing What Two Values
Cannot Write}, v1.2 (2026-09-21). Concept DOI: 10.5281/zenodo.22867336;
Version DOI: 10.5281/zenodo.22876596.

{[}3{]} Noriaki Kihara, \textbf{The Third Thought Experiment: Reading
Out the State of Two Bodies a, b from the Orbit of Two Celestial Bodies
a, b - Obtaining the Mass Ratio, the Relative Distance and the Relative
Time from Orbit and Gravitational-Wave Readouts, with Two Values of
Undetermined Meaning and the Knowledge \(G=c=M=1\) Alone}, v1.5
(2026-09-23). Concept DOI: 10.5281/zenodo.22909660; Version DOI:
10.5281/zenodo.22909661.

{[}4{]} Noriaki Kihara, \textbf{The Fourth Thought Experiment: From a
Fixed Action \(S\) to a State-Dependent Action \(S_n\) - Can the
Relativistic Gravitational Orbits Computed Externally in the Third
Thought Experiment Be Regenerated from the Sequential Interaction
\(X_{n+1}=S_nX_n\)?}, v1.0 (2026-09-23). Concept DOI:
10.5281/zenodo.22919787; Version DOI: 10.5281/zenodo.22919788.

{[}5{]} Noriaki Kihara, \textbf{The Fifth Thought Experiment: Was There
No Hidden State? - A Self-Audit of the Two-State Orbit Generation of the
Fourth Thought Experiment, Tracking the State Closure from Five States
to the Sixth Persistent State \(\mathcal N=\nu\) -}, v2.1 (2026-09-26).
Concept DOI: 10.5281/zenodo.22973086; Version DOI:
10.5281/zenodo.22973087.

\subsection{External}\label{external}

{[}6{]} P. C. Peters and J. Mathews, ``Gravitational Radiation from
Point Masses in a Keplerian Orbit,'' \emph{Physical Review}
\textbf{131}, 435-440 (1963). DOI: 10.1103/PhysRev.131.435.

{[}7{]} P. C. Peters, ``Gravitational Radiation and the Motion of Two
Point Masses,'' \emph{Physical Review} \textbf{136}, B1224-B1232 (1964).
DOI: 10.1103/PhysRev.136.B1224.

{[}8{]} L. Liu, O. Christiansen, Z.-K. Guo, R.-G. Cai, and S. P. Kim,
``Gravitational and electromagnetic radiation from binary black holes
with electric and magnetic charges: Circular orbits on a cone,''
\emph{Physical Review D} \textbf{102}, 103520 (2020). DOI:
10.1103/PhysRevD.102.103520.

{[}9{]} C. A. Benavides-Gallego and W.-B. Han, ``Gravitational Waves and
Electromagnetic Radiation from Charged Black Hole Binaries,''
\emph{Symmetry} \textbf{15}, 537 (2023). DOI: 10.3390/sym15020537.

{[}10{]} Z.-H. Zhang, T. Liu, S. Zhang, and Z.-K. Guo, ``Post-Newtonian
dynamics of charged compact binaries,'' \emph{Physical Review D}
\textbf{114}, 044088 (2026). DOI: 10.1103/4x8q-2kzm.

{[}11{]} S. Verma et al., ``Post-newtonian dynamics of radiating
charges: canonical formulation and binary inspiral laws,''
\emph{European Physical Journal C} \textbf{86}, 830 (2026). DOI:
10.1140/epjc/s10052-026-16076-2.

\begin{center}\rule{0.5\linewidth}{0.5pt}\end{center}

\section{Revision history}\label{revision-history}

\begin{itemize}
\tightlist
\item
  \textbf{v1.1 (2026-09-26):} The actual code of the initialization of
  the state generator, the local interaction rate, the sole
  \texttt{transition(z)}, the macrostep, the readout and the terminal
  interpolation was added to the main text, and the corresponding
  formulas and information flow were made explicit.
\item
  \textbf{v1.0 (2026-09-26):} First version of the sixth thought
  experiment. The charged LO quasi-circular analytic reference, the
  strict single map with eight persistent states, the five-condition
  transferability, the orbit comparison figures, the reproducibility
  information, and the scope of claims and non-claims were integrated.
\end{itemize}

\end{document}
